\documentclass[11pt]{article}

\usepackage[margin=0.92in]{geometry}
\usepackage[T1]{fontenc}
\usepackage{lmodern,microtype}
\usepackage{amsmath,amssymb,amsthm,mathtools}
\usepackage{booktabs,longtable,array,enumitem}
\usepackage[table]{xcolor}
\usepackage[most]{tcolorbox}
\usepackage[hidelinks,hypertexnames=false]{hyperref}
\usepackage[nameinlink,noabbrev]{cleveref}

\allowdisplaybreaks
\setlist[itemize]{leftmargin=1.45em,itemsep=0.2em,topsep=0.3em}
\setlength{\LTpre}{0.4em}
\setlength{\LTpost}{0.4em}

\definecolor{navy}{RGB}{25,61,98}
\definecolor{proved}{RGB}{40,118,70}
\definecolor{conditional}{RGB}{184,119,18}
\definecolor{incomplete}{RGB}{166,54,48}
\definecolor{provedbg}{RGB}{238,248,241}
\definecolor{conditionalbg}{RGB}{255,248,231}
\definecolor{incompletebg}{RGB}{253,239,238}

\hypersetup{colorlinks=true,linkcolor=navy,urlcolor=navy,
  pdftitle={Aiyagari (1994): Human-Readable Translation of the Lean Formalization},
  pdfauthor={Translation of the Lean 4 development}}

\newcommand{\R}{\mathbb R}
\newcommand{\N}{\mathbb N}
\newcommand{\E}{\mathbb E}
\newcommand{\PP}{\mathbb P}
\newcommand{\Icc}[2]{[#1,#2]}
\newcommand{\leanname}[1]{\nolinkurl{#1}}
\newcommand{\eqD}[1]{\refstepcounter{equation}\tag{D#1}}
\newcommand{\eqA}[1]{\refstepcounter{equation}\tag{A#1}}
\newcommand{\eqE}[1]{\refstepcounter{equation}\tag{E#1}}
\newcommand{\eqP}[1]{\refstepcounter{equation}\tag{P#1}}
\newcommand{\eqL}[1]{\refstepcounter{equation}\tag{L#1}}
\newcommand{\formalcheck}{\textbf{Lean status:} kernel-checked, no \texttt{sorryAx}.}
\newcommand{\missing}[1]{\par\smallskip\textbf{Missing foundational step.} #1}

\newtcolorbox{greenbox}[1]{enhanced,breakable,colback=provedbg,
  colframe=proved,boxrule=0.8pt,left=7pt,right=7pt,top=5pt,bottom=5pt,
  title={#1},fonttitle=\bfseries,coltitle=black,halign title=flush left,
  borderline west={3pt}{0pt}{proved}}
\newtcolorbox{amberbox}[1]{enhanced,breakable,colback=conditionalbg,
  colframe=conditional,boxrule=0.8pt,left=7pt,right=7pt,top=5pt,bottom=5pt,
  title={#1},fonttitle=\bfseries,coltitle=black,halign title=flush left,
  borderline west={3pt}{0pt}{conditional}}
\newtcolorbox{redbox}[1]{enhanced,breakable,colback=incompletebg,
  colframe=incomplete,boxrule=0.8pt,left=7pt,right=7pt,top=5pt,bottom=5pt,
  title={#1},fonttitle=\bfseries,coltitle=black,halign title=flush left,
  borderline west={3pt}{0pt}{incomplete}}

\newcommand{\greenstatus}{\textcolor{proved}{\textbf{GREEN -- Foundationally proved.}}}
\newcommand{\amberstatus}{\textcolor{conditional}{\textbf{AMBER -- Formally proved but incomplete.}}}
\newcommand{\redstatus}{\textcolor{incomplete}{\textbf{RED -- Intended claim not yet established.}}}

\title{\textbf{Aiyagari (1994): Human-Readable Proofs}\\[0.3em]
  \large Exact translation and adequacy audit of the Lean 4 development}
\author{Based on \texttt{i003\_replicate\_aiyagari}}
\date{}

\begin{document}
\maketitle

\begin{abstract}
This document translates the current Lean 4 formalization into conventional
mathematics.  Definitions, assumptions, and equilibrium conditions receive
the tags (D\(i\)), (A\(i\)), and (E\(i\)).  Every theorem and lemma is stated
with its actual Lean hypotheses and its proof is rewritten for human review.
All declarations are accepted by Lean without \texttt{sorry}; color instead
measures whether the declaration supplies a foundational proof of the intended
economic claim.  Green means adequate, amber means a rigorous but conditional
or narrower result, and red means that the valid Lean statement does not yet
establish the economic claim suggested by its name.
\end{abstract}

\tableofcontents

\section{Definitions}

All scalar quantities are real unless indicated otherwise.  A household state
is \(s=(a,l)\in\R^2\), and a real stochastic process on \((\Omega,\mathcal F)\)
is a map \(x:\N\times\Omega\to\R\).  An asset-supply or capital-demand curve is
a function \(\R\to\R\).

\subsection{Primitives and household accounting}

The household primitives are
\(M=(\beta,r,w,b,l_{\min},l_{\max},u,\pi_l)\), where \(\pi_l\) is a measure on
labor.  Production primitives are
\(P=(\beta,\delta,l_{\min},F,F_k,w_K,K)\).  Define
\begin{align}
 \lambda(\beta)&=\frac{1-\beta}{\beta}, \eqD{1}\label{D1}\\
 \phi^{\rm nat}(w,l_{\min},r)&=\frac{wl_{\min}}r, \eqD{2}\label{D2}\\
 \phi(b,w,l_{\min},r)&=
 \begin{cases}\min\{b,\phi^{\rm nat}(w,l_{\min},r)\},&r>0,\\b,&r\le0,
 \end{cases} \eqD{3}\label{D3}\\
 \phi^{\rm PV}(w,l_{\min},r)&=\phi^{\rm nat}(w,l_{\min},r), \eqD{4}\label{D4}\\
 \widehat a&=a+\phi, \eqD{5}\label{D5}\\
 z(w,l,r,\widehat a,\phi)&=wl+(1+r)\widehat a-r\phi, \eqD{6}\label{D6}\\
 z_{\min}(w,l_{\min},r,\phi)&=wl_{\min}-r\phi, \eqD{7}\label{D7}\\
 c_A(z)&=z-A(z). \eqD{8}\label{D8}
\end{align}
For a next-period labor realization \(l'\), resources evolve as
\begin{equation}
 T_{M,\phi}(\widehat a,l')=M.wl'+(1+M.r)\widehat a-M.r\phi.
 \eqD{9}\label{D9}
\end{equation}

\subsection{The Bellman problem}

The feasible shifted-asset set at resources \(z\) is
\begin{equation}
 \Gamma(z)=\{\widehat a:0\le \widehat a\le z\}. \eqD{10}\label{D10}
\end{equation}
For a candidate value function \(V\), set
\begin{align}
 Q_{M,\phi}V(\widehat a)
   &=\int V(T_{M,\phi}(\widehat a,l))\,\pi_l(dl), \eqD{11}\label{D11}\\
 J_{M,\phi,V}(z,\widehat a)
   &=u(z-\widehat a)+\beta Q_{M,\phi}V(\widehat a), \eqD{12}\label{D12}\\
 (\mathcal B_{M,\phi}V)(z)
   &=\sup_{\widehat a\in\Gamma(z)}J_{M,\phi,V}(z,\widehat a). \eqD{13}\label{D13}
\end{align}
A Bellman fixed point and an optimal policy mean, respectively,
\begin{align}
 \mathcal B_{M,\phi}V&=V, \eqD{14}\label{D14}\\
 A(z)&\in\Gamma(z),\qquad
 J(z,x)\le J(z,A(z))\quad(x\in\Gamma(z)). \eqD{15}\label{D15}
\end{align}

\subsection{Feasible paths and terminal conditions}

A deterministic path \((a_t,c_t,l_t)\) is feasible when
\begin{equation}
 c_t\ge0,\qquad c_t+a_{t+1}=wl_t+(1+r)a_t\quad(t\in\N). \eqD{16}\label{D16}
\end{equation}
A minimum-income feasible path imposes \(l_t=l_{\min}\) in (D16).  A stochastic
plan is feasible when (D16) holds \(\mu\)-almost surely at each date:
\begin{equation}
 \forall t,\quad c_t(\omega)\ge0,\quad
 c_t(\omega)+a_{t+1}(\omega)=wl_t(\omega)+(1+r)a_t(\omega)\quad\mu\text{-a.s.}
 \eqD{17}\label{D17}
\end{equation}
Bounded labor support and adaptedness mean
\begin{align}
 &\forall t,\quad l_{\min}\le l_t(\omega)\le l_{\max}\quad\mu\text{-a.s.},
 \eqD{18}\label{D18}\\
 &x_t:(\Omega,\mathcal F_t)\to(\R,\mathcal B(\R))\text{ is measurable for every }t.
 \eqD{19}\label{D19}
\end{align}
The pathwise and almost-sure natural debt restrictions are
\begin{align}
 &a_t\ge-\phi^{\rm nat}(w,l_{\min},r)\quad\forall t, \eqD{20}\label{D20}\\
 &\forall t,\quad a_t(\omega)\ge-\phi^{\rm nat}(w,l_{\min},r)\quad\mu\text{-a.s.}
 \eqD{21}\label{D21}
\end{align}
The Lean no-Ponzi condition requires an actual nonnegative limit:
\begin{align}
 &\exists L\ge0:\quad \frac{a_{t+1}}{(1+r)^t}\longrightarrow L, \eqD{22}\label{D22}\\
 &\text{condition (D22) holds for }\mu\text{-almost every }\omega. \eqD{23}\label{D23}
\end{align}

\subsection{Policies, kernels, and aggregation}

The auxiliary policy predicates are
\begin{align}
 \beta(1+r)e(z)&\le m(z), \eqD{24}\label{D24}\\
 z_{\min}<\widehat z,&\qquad A(z)=0\quad(z\in[z_{\min},\widehat z]), \eqD{25}\label{D25}\\
 A_0(z)&\le A_1(z),\qquad c_{A_1}(z)\le c_{A_0}(z), \eqD{26}\label{D26}\\
 z^*\le z&\Longrightarrow T(z,l)\le z. \eqD{27}\label{D27}
\end{align}
If the transition is measurable, the policy-induced Markov kernel is the law
of \(T_{M,\phi}(A(z),l')\) when \(l'\sim\pi_l\):
\begin{equation}
 \kappa_A(z,B)=\int \mathbf 1_B(T_{M,\phi}(A(z),l'))\,\pi_l(dl').
 \eqD{28}\label{D28}
\end{equation}
Invariance and stationary mean net assets are
\begin{align}
 \mu\kappa_A&=\mu, \eqD{29}\label{D29}\\
 \overline a(\phi,A,\mu)&=\int(A(z)-\phi)\,\mu(dz). \eqD{30}\label{D30}
\end{align}
For curves \(K,E_a:\R\to\R\), define
\begin{align}
 \operatorname{MC}(K,E_a,r)&:\Longleftrightarrow K(r)=E_a(r), \eqD{31}\label{D31}\\
 s(k)&=\frac{\delta k}{F(k)}, \eqD{32}\label{D32}\\
 X(r)&=K(r)-E_a(r). \eqD{33}\label{D33}
\end{align}
The full-insurance return condition is \(r=\lambda(\beta)\).  A curve is
``nonmonotone'' precisely when it is neither monotone nor antitone.  Dominance
near \(\bar r\) and one-sided explosion mean
\begin{align}
 &\exists\epsilon>0:\quad \bar r-\epsilon<r<\bar r
   \Longrightarrow E_a^{\rm certain}(r)<E_a^{\rm risky}(r), \eqD{34}\label{D34}\\
 &E_a(r)\to+\infty\ (r\uparrow\bar r),\qquad
   E_a(r)\to-\infty\ (r\downarrow\bar r). \eqD{35}\label{D35}
\end{align}
A fixed limit is slack when
\begin{equation}
 \phi^{\rm nat}(w,l_{\min},r)\le b. \eqD{36}\label{D36}
\end{equation}

\subsection{Growth, debt, and money}

The balanced-growth return is
\begin{equation}
 r_g^{\rm FI}=(1+\lambda(\beta))(1+g)^\gamma-1. \eqD{37}\label{D37}
\end{equation}
Debt-adjusted resources, net assets, and the tax-adjusted floor are
\begin{align}
 R_d(w,l,r,d,a)&=wl-rd+(1+r)a, \eqD{38}\label{D38}\\
 a^{\rm net}&=a-d, \eqD{39}\label{D39}\\
 \underline a_d&=-\frac{wl_{\min}-rd}{r}. \eqD{40}\label{D40}
\end{align}
For money \(m\), price \(p\), and a candidate sharp return \(r^\#\),
\begin{align}
 a_m&=\frac{m-1}{p}, \eqD{41}\label{D41}\\
 p^\#&=\frac{r^\#}{wl_{\min}}. \eqD{42}\label{D42}
\end{align}
Finally, \(P.\lambda=\lambda(P.\beta)\) and
\(P.s(k)=P.\delta k/P.F(k)\).

\subsection{Canonical histories and the continuous Bellman state}

The refactored household problem uses the canonical history space and product
law
\begin{equation}
 \Omega=\R^{\N},\qquad \PP=\pi_l^{\otimes\N},\qquad
 l_t(\omega)=\omega_t,\qquad \mathcal F_t=\sigma(l_0,\ldots,l_t).
 \eqD{43}\label{D43}
\end{equation}
Resources are measured above their minimum: \(x=z-z_{\min}\in\R_+\), and a
saving share \(s\in[0,1]\) represents \(\widehat a'=sz\).  Thus
\begin{equation}
 z=z_{\min}+x,\qquad c(x,s)=z(1-s),\qquad
 \Gamma_s=[0,1]. \eqD{44}\label{D44}
\end{equation}
For \(l'\in[l_{\min},l_{\max}]\), next excess resources are
\begin{equation}
 x'=\max\!\left\{0,(1+r)sz+w(l'-l_{\min})\right\}. \eqD{45}\label{D45}
\end{equation}
In the bounded branch, \(C_b(\R_+)\) has the sup norm.  For
\(V\in C_b(\R_+)\), define
\begin{align}
 Q_V(x,s)&=u(c(x,s))+\beta\int V(x'(x,s,l'))\,\pi_l(dl'),\notag\\
 (\mathbb BV)(x)&=\max_{s\in[0,1]}Q_V(x,s). \eqD{46}\label{D46}
\end{align}
The concrete \leanname{HouseholdSolution} contains the resulting value and a
maximizing policy, together with their defining equations, but contains no
uniqueness field.

For the solved bounded household, write $A^*(x)$ for the unique asset
policy and define its excess-resource transition by

\begin{equation}
 X^*(x,l)=(1+r)A^*(x)+w(l-l_{\min}).
 \eqD{47}\label{D47}
\end{equation}
If $x^*$ is a drift threshold, the compact state and its restricted
transition are
\begin{equation}
 S_{x^*}=[0,x^*],\qquad
 \widehat X^*(x,l)=X^*(x,l)\in S_{x^*}.
 \eqD{48}\label{D48}
\end{equation}
For a Markov kernel \(\kappa\), its action on laws and its Cesàro orbit
averages are
\begin{align}
 \mathcal T_\kappa(\mu)&=\kappa\circ_m\mu,
 \eqD{49}\label{D49}\\
 \overline\mu_n&=\frac1{n+1}\sum_{i=0}^{n}
   \mathcal T_\kappa^i(\mu_0).
 \eqD{50}\label{D50}
\end{align}
For a Feller kernel, the dual Markov operator and the oscillation seminorm are
\begin{align}
 (P_\kappa f)(x)&=\int f(y)\,\kappa(x,dy),
 \eqD{51}\label{D51}\\
 \operatorname{osc}(f)&=\sup_{x,y}|f(x)-f(y)|.
 \eqD{52}\label{D52}
\end{align}
The kernel-level mixing condition used below is the existence of
\(q\in[0,1)\) such that
\begin{equation}
 \operatorname{osc}(P_\kappa f)\le q\operatorname{osc}(f)
 \qquad(f\in C_b(S_{x^*})).
 \eqD{53}\label{D53}
\end{equation}
The minimum-labor and offset transitions used in the low-income-block
argument are
\begin{align}
 T_0(x)&=X^*(x,l_{\min})=(1+r)A^*(x),
 \eqD{54}\label{D54}\\
 T_\delta(x)&=(1+r)A^*(x)+w\delta,
 \qquad \delta\ge0. \eqD{55}\label{D55}
\end{align}
For an offset history $\boldsymbol\delta=(\delta_0,\delta_1,\ldots)$,
define its finite resource recursion by
\begin{equation}
 Z_0(x;\boldsymbol\delta)=x,\qquad
 Z_{n+1}(x;\boldsymbol\delta)
   =T_{\delta_n}\bigl(Z_n(x;\boldsymbol\delta)\bigr).
 \eqD{56}\label{D56}
\end{equation}
For a supported labor history, $\delta_i=l_i-l_{\min}$, so
\eqref{D56} is exactly the iterated solved transition \eqref{D47}.
The finite-step mixing condition derived below for the household kernel is
the weaker requirement that some positive iterate contracts oscillation:
\begin{equation}
 \exists m\ge1,\ \exists q\in[0,1):\qquad
 \operatorname{osc}(P_\kappa^m f)\le q\operatorname{osc}(f)
 \quad(f\in C_b(S_{x^*})).
 \eqD{57}\label{D57}
\end{equation}
For probability laws $\mu,\nu$ on $\R$, the refactored convex order and
mean-preserving-spread relation are
\begin{equation}
 \mu\preceq_{\rm cx}\nu
 \Longleftrightarrow
 \left\{\begin{array}{l}
 \int y\,d\mu=\int y\,d\nu,\\
 \int f\,d\mu\le\int f\,d\nu
 \quad\text{for every integrable convex }f.
 \end{array}\right.
 \eqD{58}\label{D58}
\end{equation}
The shifted-difference form of Miller's convex-marginal continuation class is
\begin{equation}
 V\in\mathcal M
 \Longleftrightarrow
 V\text{ is concave and }
 y\mapsto V(x_1+y)-V(x_2+y)\text{ is concave whenever }x_1\le x_2.
 \eqD{59}\label{D59}
\end{equation}
For a one-step consumption choice, define
\begin{equation}
 J_\mu(c;x,V)=u(c)+\beta\int
 V\!\left(R(x-c)+y\right)\,\mu(dy),\qquad R\ge0.
 \eqD{60}\label{D60}
\end{equation}

\section{Assumptions}

This section records exactly the structures in \texttt{Assumptions.lean}.
They are inputs, not conclusions.

The household sign and support assumptions are
\begin{equation}
 \begin{gathered}
 0<\beta<1,\quad r>0,\quad w>0,\quad b\ge0,
 \quad 0<l_{\min}<l_{\max},\\
 \pi_l(\Icc{l_{\min}}{l_{\max}})=1,\qquad
 \pi_l\text{ is a probability measure}.
 \end{gathered}
 \eqA{1}\label{A1}
\end{equation}
On \([0,\infty)\), utility is continuous, strictly increasing, strictly
concave, differentiable on \((0,\infty)\), and
\begin{equation}
 u'(c)>0\quad(c>0). \eqA{2}\label{A2}
\end{equation}
The two boundary regimes are separate:
\begin{align}
 &u\text{ is continuous at }0\text{ and has a finite right derivative there},
 \eqA{3}\label{A3}\\
 &u'(c)\longrightarrow+\infty\quad(c\downarrow0). \eqA{4}\label{A4}
\end{align}
For the bounded finite-at-zero construction, the representative utility is
twice continuously differentiable on \((0,\infty)\), has negative second
derivative there, and admits finite uniform bounds on \([0,\infty)\):
\begin{equation}
 \underline u\le u(c)\le\overline u\quad(c\ge0),\qquad u''(c)<0\quad(c>0).
 \eqA{5}\label{A5}
\end{equation}
Borrowing-limit admissibility and endpoint reachability are
\begin{align}
 z_{\min}=wl_{\min}-r\phi&\ge0, \eqA{6}\label{A6}\\
 \pi_l([l_{\min},\min\{l_{\max},l_{\min}+\epsilon\}])&>0,\qquad
 \pi_l([\max\{l_{\min},l_{\max}-\epsilon\},l_{\max}])>0
 \quad(\epsilon>0). \eqA{7}\label{A7}
\end{align}
The curvature assumptions needed by later layers are convex marginal utility
and bounded eventual relative risk aversion:
\begin{equation}
 u'\text{ is convex on }(0,\infty),\qquad
 u'\text{ is differentiable on }(c_0,\infty),\qquad
 -\frac{cu''(c)}{u'(c)}\le\overline\gamma\quad(c\ge c_0). \eqA{8}\label{A8}
\end{equation}
The positive-consumption interface assigns no economic value to \(u(0)\).  It
assumes Inada behavior and the weighted-space growth bound
\begin{equation}
 |u(c)|\le C(1+c^p),\qquad c>0,\quad C,p\ge0. \eqA{9}\label{A9}
\end{equation}
The impatience condition used only by cutoff, Euler, and drift arguments is
kept separate from the basic household signs:
\begin{equation}
 \beta(1+r)<1. \eqA{10}\label{A10}
\end{equation}
Technology assumptions are
\begin{equation}
 0<P.\beta<1,\quad P.\delta\ge0,\quad P.l_{\min}>0,
 \quad K\text{ strictly decreasing},\quad F(k)>0\ (k>0). \eqA{11}\label{A11}
\end{equation}
A resource drift certificate is
\begin{equation}
 l_{\min}\le l_{\max},\qquad z\ge z^*,\ l\in[l_{\min},l_{\max}]
 \Longrightarrow T(z,l)\le z. \eqA{12}\label{A12}
\end{equation}
The only closure input used for differentiating a labor expectation is a
local dominated-derivative certificate.  For an integrand $F(a,l)$, a
candidate derivative $F_a(a,l)$, and a reference action $a_0$, it supplies
a neighborhood $N\ni a_0$, measurability and integrability, and an
integrable envelope $g$ such that
\begin{equation}
 \partial_aF(a,l)=F_a(a,l),\qquad |F_a(a,l)|\le g(l)
 \quad(a\in N,\ \pi_l\text{-a.e. }l),\qquad \int g\,d\pi_l<\infty.
 \eqA{13}\label{A13}
\end{equation}
It does not assume the derivative of the expectation or an Euler equation;
both conclusions are derived below.
The one remaining global-policy closure used in the Proposition 4 layer is
explicitly isolated as \leanname{HighResourceConsumptionClosure}:
\begin{equation}
 c(x)\longrightarrow+\infty\qquad\text{as }x\longrightarrow+\infty.
 \eqA{14}\label{A14}
\end{equation}
This is the minimal closure requested for the global passage.  It assumes no
policy slope, affine policy bound, resource drift, absorbing interval,
stationary law, or comparative-static conclusion.  The intended later upgrade
is to derive \eqref{A14} from boundedness and concavity of the solved value and
the envelope equation.
The former Bellman-contraction, objective-concavity, law-contraction, and
invariant-law certificate structures have been removed.  Bellman contraction,
compact-Feller invariant-law existence, and the generic oscillation-mixing
theorem are now proved for concrete mathematical objects.  Finite-step
splitting is also derived below from the solved household policy, i.i.d.
lower-support reachability, and a finite-product Fubini identity.  The
remaining paper-level closures are the explicitly displayed analytic inputs
in \eqref{A13}--\eqref{A14} and positive minimum-state consumption.

\section{Equilibrium}

For a measure \(\mu\) over resources, aggregate labor, assets, and consumption
are
\begin{equation}
 L(\mu)=\int l\,d\mu,\qquad
 \mathcal A(\phi,A,\mu)=\int(A(z)-\phi)\,d\mu,\qquad
 C(A,\mu)=\int(z-A(z))\,d\mu. \eqE{1}\label{E1}
\end{equation}
A household solution consists of \((V,A)\) satisfying
\begin{equation}
 \mathcal B_{M,\phi}V=V,\qquad A\text{ satisfies (D15)},\qquad
 (z,l')\mapsto T_{M,\phi}(A(z),l')\text{ is measurable}. \eqE{2}\label{E2}
\end{equation}
Its kernel is the concrete kernel (D28).

A stationary recursive competitive equilibrium contains
\((r,w,k,C,b,\phi,M,V,A,\mu)\) satisfying the primitive matches
\begin{equation}
 M.\beta=P.\beta,\quad M.r=r,\quad M.w=w,\quad M.b=b,\quad
 M.l_{\min}=P.l_{\min}, \eqE{3}\label{E3}
\end{equation}
the household and stationary conditions
\begin{equation}
 (V,A)\text{ satisfies (E2)},\quad \mu\text{ is a probability},\quad
 \phi=\phi(b,w,P.l_{\min},r),\quad \mu\kappa_A=\mu, \eqE{4}\label{E4}
\end{equation}
firm pricing
\begin{equation}
 r=P.F_k(k)-P.\delta,\qquad w=P.w_K(k), \eqE{5}\label{E5}
\end{equation}
and clearing and aggregation
\begin{equation}
 L(\pi_l)=1,\quad \mathcal A(\phi,A,\mu)=k,\quad K(r)=k,\quad
 C(A,\mu)=C,\quad C+P.\delta k=P.F(k). \eqE{6}\label{E6}
\end{equation}
A full-insurance equilibrium satisfies
\begin{equation}
 r=P.\lambda,\quad r=P.F_k(k)-P.\delta,\quad K(r)=k,\quad
 C+P.\delta k=P.F(k). \eqE{7}\label{E7}
\end{equation}
A government-debt equilibrium has feasible paths with
\begin{equation}
 c_t+a_{t+1}=R_d(w,l_t,r,d,a_t),\qquad
 \mu\text{ probability},\quad a\text{ integrable},\quad \int a\,d\mu=d. \eqE{8}\label{E8}
\end{equation}
A monetary stationary equilibrium satisfies
\begin{equation}
 p>0,\qquad \mu\text{ probability},\quad a\text{ integrable},
 \qquad\int a\,d\mu=0. \eqE{9}\label{E9}
\end{equation}

For the complete debt-neutrality construction, translate a shifted net-asset
choice \(\widehat a\) into gross government-debt wealth by
\begin{equation}
 g_d(\widehat a)=\widehat a-\phi+d,\qquad
 \Gamma_d(z)=[d-\phi,z+d-\phi],\qquad
 c_d(z,g)=z+d-\phi-g. \eqE{10}\label{E10}
\end{equation}
The tax-adjusted transition and Bellman problem are defined directly, rather
than by assuming neutrality:
\begin{align}
 T_d(g,l')&=R_d(M.w,l',M.r,d,g)-d+\phi, \eqE{11}\label{E11}\\
 Q_d(V;z,g)&=U(c_d(z,g))+\beta\int V(T_d(g,l'))\,d\pi_l(l'), \notag\\
 (\mathcal B_dV)(z)&=\sup_{g\in\Gamma_d(z)}Q_d(V;z,g). \eqE{12}\label{E12}
\end{align}
A complete debt-free stationary equilibrium is a Bellman solution \((V,A)\)
and invariant probability \(\mu\) satisfying
\begin{equation}
 \mu\kappa_A=\mu,\qquad A-\phi\text{ integrable},\qquad
 \int(A(z)-\phi)\,d\mu(z)=0. \eqE{13}\label{E13}
\end{equation}
A complete tax-adjusted debt equilibrium requires
\begin{equation}
 \mathcal B_dV=V,\quad g_d\circ A\text{ maximizes }Q_d,\quad
 \mu\kappa_A=\mu,\quad g_d\circ A\text{ integrable},\quad
 \int g_d(A(z))\,d\mu(z)=d. \eqE{14}\label{E14}
\end{equation}

For the solved bounded household on excess resources, a full-state stationary
record \(S=(M,\phi,A,\mu)\) requires the primitive packages (A1), (A4), and
(A6), probability of \(\mu\), invariance under the actual solved kernel, and
integrability of net assets.  Its constructed mean assets are
\begin{equation}
 \mu\kappa_{M,\phi}=\mu,\qquad
 \mathfrak A(S)=\int_{\R_+}\bigl(A_{M,\phi}(x)-\phi\bigr)\,d\mu(x).
 \eqE{15}\label{E15}
\end{equation}
A solved monetary stationary equilibrium for \(M\) contains \(p>0\),
\(M.b=1/p\), and a stationary record at the effective limit
\(\phi_m=\phi(1/p,M.w,M.l_{\min},M.r)\), satisfying
\begin{equation}
 \mathfrak A(S_m)=0. \eqE{16}\label{E16}
\end{equation}
A natural-limit zero-asset equilibrium for \(M^\#\) is a stationary record at
\(\phi^{nat}(M^\#.w,M^\#.l_{\min},M^\#.r)\) satisfying
\begin{equation}
 \mathfrak A(S^\#)=0. \eqE{17}\label{E17}
\end{equation}
The primitive-matching relation
\begin{center}
\leanname{SameHouseholdEnvironmentExceptReturn}
\end{center}
requires two models to have identical \(\beta,w,l_{\min},l_{\max},U\), and
labor law; only their return and borrowing-limit bookkeeping fields may differ.

\section{Properties, lemmas, and human-readable proofs}

\subsection{Immediate equilibrium consequences}

\begin{greenbox}{Lemma E-L1: \leanname{consumption_nonneg}}
\greenstatus\quad\formalcheck

If \((V,A)\) is household equilibrium data, then \(c_A(z)\ge0\) for every
\(z\).  Indeed, optimal-policy feasibility in (D15) gives \(A(z)\le z\), so
\(c_A(z)=z-A(z)\ge0\) by (D8).
\end{greenbox}

\begin{greenbox}{Lemma E-L2: \leanname{market_clears}}
\greenstatus\quad\formalcheck

In a stationary recursive equilibrium, (E6) gives both \(K(r)=k\) and
\(\mathcal A(\phi,A,\mu)=k\).  Hence
\(K(r)=\mathcal A(\phi,A,\mu)\), which is market clearing (D31).
\end{greenbox}

\subsection{Property 1: debt limits and no Ponzi schemes}

Put \(d=\phi^{\rm nat}(w,l_{\min},r)\) and
\begin{equation}
 y_n=\frac{a_{n+1}+d}{(1+r)^n}. \eqP{1.1}\label{P1.1}
\end{equation}

\begin{greenbox}{Lemma 1.1: \leanname{minimumIncome_shifted_budget}}
\greenstatus\quad\formalcheck

For \(r\ne0\) and a minimum-income feasible path,
\begin{equation}
 c_t+(a_{t+1}+d)=(1+r)(a_t+d). \eqL{1.1}\label{L1.1}
\end{equation}
By (D2), \(rd=wl_{\min}\).  Substitute this and the budget equation (D16)
with \(l_t=l_{\min}\); rearrangement gives (L1.1).
\end{greenbox}

\begin{greenbox}{Lemma 1.2: \leanname{discountedShiftedAssets_antitone}}
\greenstatus\quad\formalcheck

If \(r>0\), the sequence (P1.1) is nonincreasing.  From (L1.1) and
\(c_{n+1}\ge0\),
\(a_{n+2}+d\le(1+r)(a_{n+1}+d)\).  Division by the positive number
\((1+r)^{n+1}\) gives \(y_{n+1}\le y_n\).
\end{greenbox}

\begin{greenbox}{Lemma 1.3: \leanname{tendsto_naturalDebtLimit_div_pow}}
\greenstatus\quad\formalcheck

For \(r>0\), \(1+r>1\), and therefore
\begin{equation}
 \frac{d}{(1+r)^n}\longrightarrow0. \eqL{1.2}\label{L1.2}
\end{equation}
This is the standard geometric-sequence limit multiplied by the constant
\(d\).
\end{greenbox}

\begin{greenbox}{Property 1: \leanname{property01_naturalDebt_iff_noPonzi}}
\greenstatus\quad\formalcheck

Let \(P=(a,c)\) be a feasible predictable plan on the canonical i.i.d.
history space.  Under bounded labor support and lower-support reachability,
\begin{equation}
 [a_t\ge-d\ \hbox{for every }t\ \hbox{a.s.}]\quad\Longleftrightarrow\quad
 \left[\exists L(\omega)\ge0:\frac{a_{t+1}(\omega)}{(1+r)^t}
 \to L(\omega)\ \hbox{a.s.}\right]. \eqP{1.2}\label{P1.2}
\end{equation}
\textit{Forward.}  Iterating and discounting the budget equation bounds the
terminal asset by current assets plus the realized present value of bounded
labor income.  The natural bound controls the terminal remainder, and
nonnegative consumption makes the discounted budget sums monotone.  Their
limit is nonnegative; countable intersection makes the conclusion hold at all
dates on one full-measure event.

\textit{Reverse.}  If the natural bound fails at date \(s+1\), predictability
gives a positive-measure event, measurable with respect to date-\(s\)
information, on which \(a_{s+1}<-d-\delta\) for some \(\delta>0\).  Choose a
long future block with every labor draw within \(\epsilon\) of \(l_{\min}\).
Independence and lower reachability give a positive-probability intersection.
On it, discounted budget summation and no Ponzi imply the opposite lower bound
once the low-income block is sufficiently long, a contradiction.
\end{greenbox}

\begin{greenbox}{Property 1 (a.s. alias): \leanname{property01_naturalDebt_iff_noPonzi_ae}}
\greenstatus\quad\formalcheck

This declaration is an explicit alias of (P1.2), retaining the same canonical
history, timing, feasibility, support, and reachability assumptions.
\end{greenbox}

\begin{greenbox}{Counterexample 1: \leanname{natural_bound_does_not_imply_noPonzi_without_feasibility}}
\greenstatus\quad\formalcheck

Let \(a_t=2^t\) for even \(t\) and \(a_t=3\cdot2^t\) for odd \(t\).  With
\(w=l_{\min}=r=1\), \(a_t\ge-1\), so (D20) holds.  Yet
\(a_{t+1}/2^t\) alternates between \(6\) and \(2\).  It is not Cauchy and has
no limit, contradicting (D22).  Thus feasibility is essential.
\end{greenbox}

\begin{greenbox}{Counterexample 2: \leanname{noPonzi_does_not_imply_natural_bound_without_feasibility}}
\greenstatus\quad\formalcheck

Let \(a_0=-2\) and \(a_t=0\) for \(t\ge1\), again with
\(w=l_{\min}=r=1\).  The discounted future path is identically zero and hence
satisfies (D22) with \(L=0\), but \(a_0=-2<-1\), violating (D20).
\end{greenbox}

\subsection{Properties 2--5: dynamic programming and policy shape}

\begin{greenbox}{Canonical i.i.d. lemmas: \leanname{canonicalLabor_map_eq},
\leanname{measure_lowerLaborBlock}, \leanname{lowerLaborBlock_pos}}
\greenstatus\quad\formalcheck

Each coordinate projection under (D43) has law \(\pi_l\).  For a finite index
set \(I=\{t,\ldots,t+n\}\) and
\(B_\epsilon=[l_{\min},\min\{l_{\max},l_{\min}+\epsilon\}]\), the product
measure cylinder formula gives
\begin{equation}
 \PP\{l_i\in B_\epsilon\ \forall i\in I\}
   =\prod_{i\in I}\pi_l(B_\epsilon)>0. \eqL{2.2}\label{L2.2}
\end{equation}
The equality is independence of the canonical coordinates; strict positivity
uses lower reachability (A7), factor by factor.  This is the probability
input used in the reverse direction of Property 1.
\end{greenbox}

\begin{greenbox}{Compact maximization lemmas:
\leanname{exists_compactMaximum_eq}, \leanname{compactMaximum_le},
\leanname{abs_compactMaximum_sub_le},
\leanname{compactMaximum_lipschitz},
\leanname{continuous_compactMaximum},
\leanname{maximizeObjective_norm_sub_le},
\leanname{maximizeObjective_lipschitz},
\leanname{continuous_maximizeObjective}}
\greenstatus\quad\formalcheck

If \(f,g\) are continuous on a nonempty compact action space, both maxima are
attained.  Evaluating \(g\) at a maximizer of \(f\), and conversely, gives
\begin{equation}
 |\max f-\max g|\le\|f-g\|_\infty. \eqL{2.3}\label{L2.3}
\end{equation}
Applying this state by state proves that maximization of a bounded continuous
state-action objective is nonexpansive and continuous.
\end{greenbox}

\begin{greenbox}{Expectation lemmas: \leanname{abs_integral_le_uniform_bound},
\leanname{abs_integral_bcf_sub_le},
\leanname{expectedContinuationBCF_norm_sub_le},
\leanname{expectedContinuationBCF_lipschitz},
\leanname{continuous_expectedContinuationBCF}}
\greenstatus\quad\formalcheck

For a probability measure and bounded continuous \(V,W\),
\begin{equation}
 \left|\int(V-W)\,d\pi_l\right|
 \le\int|V-W|\,d\pi_l\le\|V-W\|_\infty. \eqL{2.4}\label{L2.4}
\end{equation}
The parametric-integral theorem and continuity of (D45) show that expected
continuation is bounded continuous.  Taking the supremum proves
nonexpansiveness of the continuation operator.
\end{greenbox}

\begin{greenbox}{Solved bounded Bellman core:
\leanname{continuousBellmanOperator_norm_sub_le},
\leanname{continuousBellmanOperator_contracting},
\leanname{continuous_continuousBellmanOperator_reward_value}}
\greenstatus\quad\formalcheck

Combining (L2.3)--(L2.4) with (D46) yields directly
\begin{equation}
 \|\mathbb BV-\mathbb BW\|_\infty
 \le\beta\|V-W\|_\infty. \eqL{2.5}\label{L2.5}
\end{equation}
Since \(0\le\beta<1\), this is a contraction on the complete space
\(C_b(\R_+)\).  No contraction certificate is supplied as a hypothesis.
\end{greenbox}

\begin{greenbox}{Parameterized Banach fixed points:
\leanname{parameterizedContractionFixedPoint_isFixedPt},
\leanname{dist_parameterizedContractionFixedPoint_le},
\leanname{continuous_parameterizedContractionFixedPoint},
\leanname{continuous_boundedBellmanValue_reward}}
\greenstatus\quad\formalcheck

Let $T_\theta:X\to X$ be a family of maps on a complete metric space, all
contracting with the same coefficient $K<1$, and let $V_\theta$ denote the
Banach fixed point.  Applying the contraction residual estimate for
$T_\theta$ at the old fixed point $V_\eta$ gives
\begin{equation}
 d(V_\theta,V_\eta)
 \le \frac{d\bigl(V_\eta,T_\theta V_\eta\bigr)}{1-K}.
 \eqL{2.5a}\label{L2.5a}
\end{equation}
If $(\theta,V)\mapsto T_\theta V$ is jointly continuous and
$\theta_n\to\theta$, then the numerator in \eqref{L2.5a}, with
$\eta=\theta$, converges to
$d(V_\theta,T_\theta V_\theta)=0$.  The squeeze theorem therefore proves
$V_{\theta_n}\to V_\theta$.  This establishes parameter continuity of
fixed points without compactness of the value-function space or a supplied
fixed-point-continuity field.

For the concrete bounded Bellman operator, the expectation map is
nonexpansive, maximization is nonexpansive, and continuous addition and scalar
multiplication show that $(R,V)\mapsto\mathbb B_RV$ is jointly continuous
when the transition and discount factor are fixed.  Hence
$R\mapsto V_R$ is now proved continuous.  Extending this calculation to the
price-dependent transition is the remaining model-specific part.
\end{greenbox}

\begin{greenbox}{Compact parameterized Bellman and policy theorem:
\leanname{continuous_boundedContinuousFamily},
\leanname{continuous_compactParameterizedExpectation},
\leanname{continuous_compactParameterizedBellmanOperator},
\leanname{continuous_compactParameterizedBellmanValue},
\leanname{continuous_continuousFamilyArgmax},
\leanname{continuous_compactParameterizedSolvedPolicy}}
\greenstatus\quad\formalcheck

Let $\Theta$ be locally compact and metrizable and let the state, action, and
shock spaces $X,S,L$ be common compact metric spaces.  Suppose the primitive
reward $R_\theta\in C_b(X\times S)$ is continuous in the sup norm and the
transition $g(\theta,x,s,l)$ is jointly continuous.  For fixed $V$, the
Heine--Cantor theorem upgrades joint continuity to uniform continuity over
the compact fiber.  Probability integration then gives
\begin{equation}
 \theta_n\to\theta
 \quad\Longrightarrow\quad
 \sup_{x,s}\left|\int V(g(\theta_n,x,s,l))\,d\pi_l
 -\int V(g(\theta,x,s,l))\,d\pi_l\right|\to0. \eqL{2.5b}\label{L2.5b}
\end{equation}
Together with the one-Lipschitz estimate in $V$, this proves joint continuity
of expected continuation in $(\theta,V)$.  Hence
\begin{equation}
 (\theta,V)\longmapsto
 \max_{s\in S}\left\{R_\theta(x,s)+
 \beta\int V(g(\theta,x,s,l))\,d\pi_l\right\} \eqL{2.5c}\label{L2.5c}
\end{equation}
is jointly continuous as a map into $C_b(X)$ and is a uniform
$\beta$-contraction in $V$.  Equation \eqref{L2.5a} therefore proves sup-norm
continuity of the solved value $V_\theta$.

The solved objective is consequently jointly continuous in $(\theta,x,s)$.
If primitive strict concavity makes its maximizer unique, every subsequential
limit of maximizers remains optimal by continuity and must equal that unique
maximizer.  Compactness then proves joint continuity of the policy in
$(\theta,x)$.  This is a proved maximum-theorem conclusion, not a supplied
policy-shape or policy-continuity hypothesis.

\textbf{Scope.}  This theorem requires one common compact state space that is
preserved by every feasible transition in the Bellman problem.  Property 6
currently gives an absorbing interval only for the \emph{solved policy}; it
does not show that every feasible Bellman action stays in that interval.
Thus applying the compact theorem to the untruncated global Aiyagari problem
still requires either a uniform truncation-equivalence theorem or a weighted
global function space.  The generic theorem itself is foundationally proved.
\end{greenbox}

\begin{greenbox}{Value and optimizer construction:
\leanname{boundedBellmanValue_isFixedPoint},
\leanname{boundedBellmanValue_unique},
\leanname{boundedAiyagariValue_isFixedPoint},
\leanname{boundedAiyagariValue_unique},
\leanname{boundedAiyagariPolicy_isMax}}
\greenstatus\quad\formalcheck

Banach's theorem applied to (L2.5) constructs a unique
\(V^*\in C_b(\R_+)\) satisfying \(\mathbb BV^*=V^*\).  For each \(x\),
continuity of \(s\mapsto Q_{V^*}(x,s)\) and compactness of \([0,1]\) produce
a maximizing share \(s^*(x)\).  Hence
\begin{equation}
 Q_{V^*}(x,s)\le Q_{V^*}(x,s^*(x))\quad
 (x\ge0,\ s\in[0,1]). \eqL{2.6}\label{L2.6}
\end{equation}
The later strict-concavity and maximum-theorem lemmas upgrade this initial
selector to the unique continuous shifted-asset policy used in Property 2.
\end{greenbox}

\begin{greenbox}{Aiyagari feasibility and transition:
\leanname{shareConsumption_nonneg},
\leanname{continuous_shareNextExcess}}
\greenstatus\quad\formalcheck

Under (A6), \(z=z_{\min}+x\ge0\); since \(s\in[0,1]\), (D44) gives
\(c=z(1-s)\ge0\).  Formula (D45) is continuous because it is assembled from
continuous arithmetic operations and \(y\mapsto\max\{0,y\}\); its codomain is
\(\R_+\) because the maximum is nonnegative.
\end{greenbox}

\begin{greenbox}{Lemma 2.1: \leanname{beta_mul_one_add_timePreferenceRate}}
\greenstatus\quad\formalcheck

For \(\beta\ne0\), direct substitution in (D1) gives
\begin{equation}
 \beta(1+\lambda(\beta))=1. \eqL{2.7}\label{L2.7}
\end{equation}
\end{greenbox}

\begin{greenbox}{Lemma 2.2: \leanname{boundedAiyagariValue_unique}}
\greenstatus\quad\formalcheck

The concrete bounded Bellman map is a \(\beta\)-contraction in the ordinary
sup norm.  Thus, if \(V,W\) are fixed points,
\[
 \|V-W\|_\infty=\|\mathcal BV-\mathcal BW\|_\infty
 \le\beta\|V-W\|_\infty.
\]
Since \(0<\beta<1\), the distance is zero and \(V=W\).  Banach's theorem also
constructs the fixed point from the primitive Bellman operator.
\end{greenbox}

\begin{greenbox}{Lemma 2.3: \leanname{boundedOptimalAsset_unique}}
\greenstatus\quad\formalcheck

Bellman iteration preserves concavity, so the solved continuation value is
concave.  Adding strictly concave current utility makes the objective strictly
concave in shifted assets on \([0,z]\).  Two distinct maximizers would violate
the defining strict-concavity inequality at their midpoint; hence the
optimizer is unique.
\end{greenbox}

\begin{greenbox}{Envelope theorem: \leanname{boundedAiyagari_envelope_of_consumption_pos}}
\greenstatus\quad\formalcheck

Fix $x>0$, write $a=A(x)$, $z=z_{\min}+x$, and suppose
$c=z-a>0$.  Holding $a$ fixed as current resources vary gives
\[
 s(y)=u(z_{\min}+y-a)+\beta G(a),\qquad
 G(a)=\int V((1+r)a+w(l-l_{\min}))\,d\pi_l.
\]
This action remains feasible for $y$ near $x$, hence $s(y)\le V(y)$,
while Bellman optimality gives $s(x)=V(x)$.  Concavity supplies both
one-sided derivatives of $V$.  The touching lower support and the ordering
of concave one-sided slopes force both to equal $s'(x)=u'(c)$.  Thus
\begin{equation}
 V'(x)=u'(c(x)). \eqP{2.1}\label{P2.1}
\end{equation}
No derivative of the policy is used or assumed.
\end{greenbox}

\begin{greenbox}{Euler system: \leanname{boundedAiyagari_euler_inequality}}
\greenstatus\quad\formalcheck

Under \eqref{A13}, dominated differentiation first proves
\[
 G'(a)=\int (1+r)V'((1+r)a+w(l-l_{\min}))\,d\pi_l.
\]
The objective derivative is $-u'(z-a)+\beta G'(a)$.  At the lower
asset boundary, every sufficiently small rightward deviation is feasible and
cannot improve the objective, so this derivative is nonpositive.  If $a>0$,
positive consumption also gives $a<z$, making the optimizer interior;
Fermat's theorem then makes the derivative zero.  Therefore
\begin{equation}
 u'(c(x))\ge \beta(1+r)\int V'(x'(x,l))\,d\pi_l,
 \quad\text{with equality if }A(x)>0. \eqP{2.2}\label{P2.2}
\end{equation}
\end{greenbox}

\begin{amberbox}{Property 2: \leanname{property02_value_and_policy_regularity}}
\amberstatus\quad\formalcheck

Under the bounded finite-at-zero assumptions, Lean constructs
\leanname{solveBoundedHousehold}, denoted \(S\), from primitives alone.  The
theorem proves: \(S.V\) is the unique bounded continuous fixed point; \(S.V\)
is concave; each feasible objective is strictly concave; \(S.A\) is feasible,
optimal, unique, and continuous.  Policy continuity is proved by the maximum
theorem: every convergent state sequence has a convergent subsequence of
saving shares, continuity passes optimality to the limit, and optimizer
uniqueness identifies that limit with \(S.A(x)\).  The numbered declaration
now also exposes the envelope identity \eqref{P2.1} and, whenever the explicit
local domination certificate \eqref{A13} is supplied, the Euler inequality
and complementary-slackness equality \eqref{P2.2}.

\missing{Prove the domination certificate globally from primitive regularity.
For parameter continuity, either prove a common compact truncation is exactly
equivalent to the unbounded household problem on admissible parameter sets,
then apply \eqref{L2.5b}--\eqref{L2.5c}, or construct the global weighted
function space directly.  Finally specialize the now-proved compact unique
argmax theorem to the household objective and construct the distinct
positive-consumption weighted Inada theory.}
\end{amberbox}

\begin{amberbox}{Property 3: \leanname{property03_binding_region_at_low_resources}}
\amberstatus\quad\formalcheck

Assume impatience \(\beta(1+r)<1\), positive optimal consumption at the
minimum excess-resource state, and the local dominated-differentiation
condition \eqref{A13} at every policy choice.  Put
\(m(x)=u'(c(x))\).  The envelope and Euler equations imply that whenever
\(A(x)>0\),
\begin{equation}
 m(x)=\beta(1+r)\int m(x'(A(x),l))\,d\pi_l(l)
 \leq \beta(1+r)m(0). \eqL{3.1}\label{L3.1}
\end{equation}
Indeed, \(A(x)>0\) makes every next state positive; monotonicity of
consumption and concavity of utility give
\(m(x'(A(x),l))\leq m(0)\).  At \(x=0\), (L3.1) and
\(\beta(1+r)<1\) contradict \(m(0)>0\), proving \(A(0)=0\).

Continuity of consumption and of marginal utility at the positive number
\(c(0)\) yields a \(\delta>0\) such that, for \(0\leq x<\delta\),
\[
 \beta(1+r)m(0)<m(x).
\]
If saving were positive at such an \(x\), (L3.1) would give the reverse weak
inequality.  Taking \(\widehat x=\delta/2\), the Lean theorem therefore
constructs a genuinely nondegenerate cutoff:
\begin{equation}
 \widehat x>0,\qquad
 0\leq x\leq\widehat x\Longrightarrow
 A(x)=0\quad\hbox{and}\quad c(x)=z_{\min}+x.
 \eqP{3.1}\label{P3.1}
\end{equation}
The supporting declarations are
\leanname{boundedAiyagariAssetPolicy_zero_at_minimum},
\leanname{boundedAiyagari_marginal_le_minimum_of_policy_pos}, and
\leanname{boundedAiyagari_binding_cutoff}.

\missing{Derive positive minimum-state consumption and the global family of
dominated-differentiation certificates from the primitive finite-at-zero
utility and labor assumptions, rather than supplying them to the numbered
theorem.}
\end{amberbox}

\begin{redbox}{Inada claim: \leanname{property03_inada_exception}}
\redstatus\quad\formalcheck

The Lean statement assumes \(A(z)>0\) for every \(z>z_{\min}\) and concludes
\(A(z)\ne0\).  The proof is the order fact ``strictly positive implies
nonzero.''

\missing{Use (A4), feasibility, and optimality to prove policy positivity.
The current theorem neither mentions nor uses the Inada assumption.}
\end{redbox}

\begin{amberbox}{Property 4: \leanname{property04_monotonicity_and_one_for_one_response}}
\amberstatus\quad\formalcheck

For the actual solved bounded household, Lean now proves that both assets and
consumption are weakly increasing and continuous.  Revealed preference and
strict concavity rule out a decrease in assets.  A second revealed-preference
argument exchanges the two consumption choices; strict optimality then
contradicts concavity of expected continuation value if consumption falls.
Consequently, for \(x_1\le x_2\),
\begin{equation}
 0\le A(x_2)-A(x_1)\le x_2-x_1. \eqP{4.1}\label{P4.1}
\end{equation}
Thus the weak derivative bounds \(0\le A'\le1\) follow at every point where
the policy is differentiable, without assuming them as policy-shape inputs.

\missing{Use the envelope/Euler system plus explicit smoothness and
nondegeneracy to prove strict monotonicity and \(0<A'(z)<1\) off the binding
region; combine this with Property 3 for the one-for-one binding response.}
\end{amberbox}

\begin{greenbox}{Lemma 5.1: \leanname{consumptionShiftedDown_of_policyShiftedUp}}
\greenstatus\quad\formalcheck

If \(A_0(z)\le A_1(z)\), then by (D8)
\(c_{A_1}(z)=z-A_1(z)\le z-A_0(z)=c_{A_0}(z)\).
\end{greenbox}

\begin{greenbox}{Same-law aggregation:
\leanname{stationaryMeanAssets_mono_same_law}}
\greenstatus\quad\formalcheck

If both policies are integrated against the \emph{same} probability law
$\mu$ and $A_0(z)\le A_1(z)$ pointwise, integrability and monotonicity of the
integral give
\[
 \int (A_0(z)-\phi)\,d\mu(z)
 \le \int (A_1(z)-\phi)\,d\mu(z).
\]
This lemma deliberately fixes the law; it does not compare the two endogenous
invariant laws generated by different shock distributions.
\end{greenbox}

\begin{greenbox}{Convex-order and saving-gap comparison:
\leanname{integral_mono_of_convexOrder_concave},
\leanname{continuationSavingGap_mono_of_meanPreservingSpread},
\leanname{oneStepObjectiveSavingGap_mono_of_meanPreservingSpread}}
\greenstatus\quad\formalcheck

If $\mu\preceq_{\rm cx}\nu$, applying \eqref{D58} to the negative of an
integrable concave function gives
\begin{equation}
 f\text{ concave}\quad\Longrightarrow\quad
 \int f\,d\nu\le\int f\,d\mu.
 \eqL{5.1}\label{L5.1}
\end{equation}
Let $c_L\le c_H$.  Since $R(x-c_H)\le R(x-c_L)$ and $V\in\mathcal M$,
\[
 y\longmapsto
 V(R(x-c_H)+y)-V(R(x-c_L)+y)
\]
is concave.  Applying \eqref{L5.1} and rearranging proves Miller's
four-point comparison
\begin{equation}
 \begin{split}
 &\int V(R(x-c_L)+y)\,d\mu-\int V(R(x-c_H)+y)\,d\mu\\
 &\qquad\le
 \int V(R(x-c_L)+y)\,d\nu-\int V(R(x-c_H)+y)\,d\nu.
 \end{split}
 \eqL{5.2}\label{L5.2}
\end{equation}
Multiplication by $\beta\ge0$ and cancellation of current utility gives the
same inequality for $J_\mu(c_L)-J_\mu(c_H)$ and
$J_\nu(c_L)-J_\nu(c_H)$.
\end{greenbox}

\begin{greenbox}{One-step optimizer comparison:
\leanname{optimalConsumption_antitone_of_meanPreservingSpread_sameContinuation}}
\greenstatus\quad\formalcheck

Let $c_0$ uniquely maximize $J_\mu$ and $c_1$ uniquely maximize $J_\nu$,
with the same continuation function $V\in\mathcal M$.  If $c_0<c_1$,
baseline optimality makes $J_\mu(c_0)-J_\mu(c_1)\ge0$.  Equation
\eqref{L5.2} then gives $J_\nu(c_0)-J_\nu(c_1)\ge0$.  Risky-law optimality
gives the reverse inequality, so both choices maximize $J_\nu$; uniqueness
forces $c_0=c_1$, a contradiction.  Therefore
\begin{equation}
 c_1\le c_0,\qquad x-c_0\le x-c_1.
 \eqL{5.3}\label{L5.3}
\end{equation}
\end{greenbox}

\begin{amberbox}{Property 5: \leanname{property05_mean_preserving_spread}}
\amberstatus\quad\formalcheck

The numbered theorem now states \eqref{L5.3} for arbitrary continuous shock
laws related by the genuine convex order \eqref{D58}; it no longer uses
arbitrary policies or unrelated Dirac laws.  Thus the one-step
Sibley--Miller single-crossing core is foundationally proved.

\missing{In the infinite-horizon household, the two shock laws generate
different fixed-point continuation values $V_\mu$ and $V_\nu$.  Complete
Miller's Lemmas 1--3 in the present continuous-state Bellman space: prove
that the Bellman operator preserves \eqref{D59}, that its fixed point remains
in this class, and that $V_\nu-V_\mu$ is nondecreasing.  This will let
\eqref{L5.3} compare the two solved policies rather than two objectives with
a common continuation value.  Aggregate stationary assets must then be
handled under the two endogenous invariant laws; pointwise ordering alone
does not order those two integrals.}
\end{amberbox}

\subsection{Properties 6--10: stationary dynamics and asset supply}

\begin{greenbox}{RRA growth lemmas:
\leanname{EventualRelativeRiskAversionBound.marginal_ratio_le_rpow},
\leanname{EventualRelativeRiskAversionBound.tendsto_marginal_ratio_add_const}}
\greenstatus\quad\formalcheck

For $c_0\le a\le b$, differentiate
$g(c)=\log u'(c)+\overline\gamma\log c$.  Assumption \eqref{A8}
gives
\[
 g'(c)=\frac{u''(c)}{u'(c)}+\frac{\overline\gamma}{c}\ge0.
\]
The mean-value monotonicity theorem and exponentiation therefore yield
\begin{equation}
 \frac{u'(a)}{u'(b)}\le
 \left(\frac ba\right)^{\overline\gamma}.
 \eqL{6.1}\label{L6.1}
\end{equation}
For each fixed $d\ge0$, decreasing marginal utility supplies the lower bound
$1\le u'(c)/u'(c+d)$, while \eqref{L6.1} supplies the upper bound
$(1+d/c)^{\overline\gamma}$.  Since the latter converges to one,
\begin{equation}
 \lim_{c\to\infty}\frac{u'(c)}{u'(c+d)}=1.
 \eqL{6.2}\label{L6.2}
\end{equation}
These are the relative-risk-aversion steps in Aiyagari's 1993 Proposition 4;
no drift or affine policy bound is assumed in either lemma.
\end{greenbox}
\begin{greenbox}{Bounded labor-gap comparison:
\leanname{rra_marginal_ratio_bounded_gap_tendsto_one}}
\greenstatus\quad\formalcheck

Let $c$ be increasing, $x_-(x)\le x_+(x)$, and suppose
$c(x_+(x))-c(x_-(x))\le d$ for a fixed $d\ge0$.  If
$c(x_-(x))\to\infty$, marginal-utility monotonicity and \eqref{L6.2} give
\begin{equation}
 1\le \frac{u'(c(x_-(x)))}{u'(c(x_+(x)))}
 \le \frac{u'(c(x_-(x)))}{u'(c(x_-(x))+d)}\longrightarrow1.
 \eqL{6.3}\label{L6.3}
\end{equation}
The squeeze theorem proves that the middle ratio converges to one.
\end{greenbox}

\begin{greenbox}{Proposition 4, unbounded-policy branch:
\leanname{bounded_resource_dynamics_of_policy_tendsto_atTop}}
\greenstatus\quad\formalcheck

Write $R=1+r$, let $A$ and $c$ be the solved Bellman policy and its induced
consumption, and assume $A(x)\to\infty$.  The lowest and highest possible next
excess-resource states are
\begin{equation}
 x_-(x)=R A(x),\qquad
 x_+(x)=R A(x)+w(l_{\max}-l_{\min}).
 \eqP{6.1}\label{P6.1}
\end{equation}
Policy and consumption monotonicity imply
$x_-(x)\le x_+(x)$ and
$0\le c(x_+(x))-c(x_-(x))\le w(l_{\max}-l_{\min})$.
By \eqref{A14}, $c(x_-(x))\to\infty$; hence \eqref{L6.3} applies.

Choose $\varepsilon>0$ so that
\begin{equation}
 \beta R(1+\varepsilon)<1. \eqP{6.2}\label{P6.2}
\end{equation}
For all sufficiently large $x$, \eqref{L6.3} yields
$u'(c(x_-))<(1+\varepsilon)u'(c(x_+))$.  Every labor realization produces
$X'=RA(x)+w(l-l_{\min})\in[x_-,x_+]$.  The envelope equation and decreasing
marginal utility therefore imply
\[
 V'(X')=u'(c(X'))\le u'(c(x_-))
 <(1+\varepsilon)u'(c(x_+)).
\]
Integrating this inequality and using the interior Euler equality gives
\begin{equation}
 V'(x)=u'(c(x))
 \le\beta R(1+\varepsilon)u'(c(x_+))<V'(x_+).
 \eqP{6.3}\label{P6.3}
\end{equation}
If $x<x_+$, concavity of $V$ instead forces
$V'(x_+)\le [V(x_+)-V(x)]/(x_+-x)\le V'(x)$, contradicting
\eqref{P6.3}.  Thus $x_+(x)\le x$ for all sufficiently large $x$, and so
every feasible next state satisfies $X'\le x$.
\end{greenbox}

\begin{amberbox}{Property 6: \leanname{property06_bounded_resource_dynamics}}
\amberstatus\quad\formalcheck

The theorem divides into exhaustive cases for the solved policy.  If
$A$ is bounded above by $K$, then
\[
 X'\le RK+w(l_{\max}-l_{\min}),
\]
so any larger excess-resource state is a drift threshold.  If $A$ is not
bounded above, its proved monotonicity implies $A(x)\to\infty$, and the
preceding unbounded-policy theorem applies.  Consequently
\begin{equation}
 \exists x^*<\infty\ \forall x\ge x^*\ \forall l\in[l_{\min},l_{\max}],
 \qquad X'(x,l)\le x. \eqP{6.4}\label{P6.4}
\end{equation}

\missing{The drift conclusion is now proved from the solved household policy,
the Euler and envelope equations, bounded shocks, impatience, and the primitive
RRA bound.  To upgrade the paper-level status to green, derive the sole
remaining global closure \eqref{A14} from the bounded Bellman solution rather
than assuming it.}
\end{amberbox}

\begin{greenbox}{Feller and monotonicity foundations:
\leanname{iidTransitionFeller_of_continuous},
\leanname{feller_solvedResourceTransition},
\leanname{monotone_solvedResourceTransition}}
\greenstatus\quad\formalcheck

The solved transition on excess resources is
\[
 X'(x,l')=(1+r)A(x)+w(l'-l_{\min}).
\]
Policy continuity and ordinary algebra make \((x,l')\mapsto X'(x,l')\)
jointly continuous.  For every bounded continuous test function \(f\), the
integrand \(f(X'(x,l'))\) is jointly continuous and uniformly bounded on the
compact labor support.  Continuity of parametric integration therefore gives
\(x\mapsto\int f(X'(x,l'))\,\pi_l(dl')\) continuous: the household kernel is
Feller.  Policy monotonicity and \(1+r>0\) also give
\(x_1\le x_2\Rightarrow X'(x_1,l')\le X'(x_2,l')\) for every shock.
\end{greenbox}

\begin{greenbox}{Absorbing compact household kernel:
\leanname{solvedNextExcess_le_driftThreshold},
\leanname{feller_absorbingSolvedKernel},
\leanname{exists_invariantProbability_absorbingSolvedKernel}}
\greenstatus\quad\formalcheck

Let \(x^*\) satisfy the Property 6 drift inequality.  If \(x\le x^*\),
monotonicity of \eqref{D47} gives
\[
 X^*(x,l)\le X^*(x^*,l)\le x^*.
\]
Thus \(S_{x^*}\) in \eqref{D48} is positively invariant.  Restricting the
jointly continuous transition to this interval preserves joint continuity,
and hence the induced kernel on \(S_{x^*}\) is Feller.
\end{greenbox}

\begin{greenbox}{Krylov--Bogoliubov existence:
\leanname{cesaro_residual_tendsto_zero_of_map_average},
\leanname{exists_fixedPoint_of_compact_cesaro},
\leanname{continuous_markovLawEvolution_of_feller},
\leanname{markovLawEvolution_probabilityCesaroAverage},
\leanname{integral_iterate_markovLawEvolution},
\leanname{exists_invariantProbability_of_feller_compact},
\leanname{solvedHousehold_invariantProbability_exists}}
\greenstatus\quad\formalcheck

Kernel linearity maps \eqref{D50} into the shifted orbit average.  Therefore,
for every bounded continuous \(f\), telescoping gives
\begin{equation}
 \left|\int f\,d\mathcal T_\kappa(\overline\mu_n)
       -\int f\,d\overline\mu_n\right|
 \le \frac{2\lVert f\rVert_\infty}{n+1}\longrightarrow0.
 \eqL{7.1}\label{L7.1}
\end{equation}
The space of probability laws on compact \(S_{x^*}\) is compact.  Hence a
subsequence \(\overline\mu_{n_j}\) converges weakly to some \(\mu^*\).
Feller continuity makes \(\mathcal T_\kappa\) continuous, while
\eqref{L7.1} forces both \(\mathcal T_\kappa(\overline\mu_{n_j})\) and
\(\overline\mu_{n_j}\) to have the same limit.  Uniqueness of weak limits
then yields \(\mathcal T_\kappa(\mu^*)=\mu^*\).  Combining this theorem with
Property 6 proves existence for the actual solved household kernel, without
assuming an invariant law or a contraction of laws.
\end{greenbox}

\begin{greenbox}{Oscillation mixing, uniqueness, and stability:
\leanname{bcfOscillation_iterate_le},
\leanname{bcfOscillation_fellerMarkovOperator_le_two},
\leanname{bcfOscillation_mul_iterate_le},
\leanname{exists_unique_invariant_and_tendsto_of_feller_mixing}}
\greenstatus\quad\formalcheck

Under \eqref{D53}, induction gives
\begin{equation}
 \operatorname{osc}(P_\kappa^n f)
 \le q^n\operatorname{osc}(f).
 \eqL{7.2}\label{L7.2}
\end{equation}
For any probability laws \(\mu,\nu\), centering \(f\) at an arbitrary state
and applying the sup-norm integral bound gives
\begin{equation}
 \left|\int f\,d\mu-\int f\,d\nu\right|
 \le2\operatorname{osc}(f).
 \eqL{7.3}\label{L7.3}
\end{equation}
Duality of \eqref{D49} and \eqref{D51} yields
\(\int f\,d\mathcal T_\kappa^n\mu=\int P_\kappa^n f\,d\mu\).
Applying \eqref{L7.3} to \(P_\kappa^n f\), then \eqref{L7.2}, proves
\begin{equation}
 \left|\int f\,d\mathcal T_\kappa^n\mu_0-
       \int f\,d\mu^*\right|
 \le2q^n\operatorname{osc}(f)\longrightarrow0.
 \eqL{7.4}\label{L7.4}
\end{equation}
Thus every initial law converges weakly to an invariant law.  Applying the
same result to a second invariant law makes a constant sequence converge to
\(\mu^*\), proving uniqueness.
\end{greenbox}

\begin{greenbox}{Finite-step mixing controls the original chain:
\leanname{tendsto_bcfOscillation_iterate_of_finiteStepMixing},
\leanname{tendsto_iterate_of_kernelIterateOscillationMixing},
\leanname{exists_unique_invariant_and_tendsto_of_feller_finiteStepMixing}}
\greenstatus\quad\formalcheck

Assume \eqref{D57}.  Write $n=k_nm+j_n$, with
$k_n=\lfloor n/m\rfloor$ and $0\le j_n<m$.  A general Markov operator
expands oscillation by at most a factor two, whereas each complete block
contracts it by $q$.  Hence
\begin{equation}
 \operatorname{osc}(P_\kappa^n f)
 \le 2^{j_n}q^{k_n}\operatorname{osc}(f)
 \le 2^m q^{\lfloor n/m\rfloor}\operatorname{osc}(f)
 \longrightarrow0.
 \eqL{7.9}\label{L7.9}
\end{equation}
Substituting \eqref{L7.9} for \eqref{L7.2} in the duality argument
\eqref{L7.3}--\eqref{L7.4} proves convergence at every original date, not
only along multiples of $m$.  Invariance then proves uniqueness exactly as
before.  Thus no separate skeleton-to-original-chain assumption is needed.
\end{greenbox}

\begin{greenbox}{Common-shock splitting implies mixing:
\leanname{kernelOscillationMixing_of_oneStepSplitting}}
\greenstatus\quad\formalcheck

Suppose a measurable shock event \(S\) has probability \(p>0\) and, on
\(S\), the transition is independent of the current state.  For any
\(x,y\), the integrands defining \(P_\kappa f(x)\) and
\(P_\kappa f(y)\) coincide on \(S\).  On \(S^c\) their absolute difference
is at most \(\operatorname{osc}(f)\).  Hence
\begin{equation}
 |P_\kappa f(x)-P_\kappa f(y)|
 \le (1-p)\operatorname{osc}(f),
 \eqL{7.5}\label{L7.5}
\end{equation}
so \eqref{D53} holds with \(q=1-p<1\).
\end{greenbox}

\begin{greenbox}{Minimum-income strict descent:
\leanname{minimumLaborTransition_strictly_decreases}}
\greenstatus\quad\formalcheck

Assume $x>0$.  If $A^*(x)=0$, then \eqref{D54} gives $T_0(x)=0<x$.
Otherwise the asset choice is interior and the Euler equality applies.
Suppose toward a contradiction that $T_0(x)\ge x$.  Since every supported
labor realization satisfies $l'\ge l_{\min}$,
$X^*(x,l')\ge T_0(x)\ge x$.  Monotonicity of solved consumption and
decreasing marginal utility therefore imply, for every $l'$, that
\[
 V'\bigl(X^*(x,l')\bigr)
 =u'\bigl(c^*(X^*(x,l'))\bigr)\le u'(c^*(x)).
\]
Substitution into the interior Euler equation yields
\begin{equation}
 u'(c^*(x))
 \le \beta(1+r)u'(c^*(x))<u'(c^*(x)),
 \eqL{7.6}\label{L7.6}
\end{equation}
where the final inequality uses \eqref{A9} and positive marginal utility.
This contradiction proves $T_0(x)<x$ at every $x>0$.
\end{greenbox}

\begin{greenbox}{Uniform finite entrance under minimum labor:
\leanname{tendsto_iterate_zero_of_monotone_strictDecrease},
\leanname{exists_uniform_iterate_lt_of_monotone_strictDecrease},
\leanname{minimumLaborBlock_uniformly_enters_bindingRegion}}
\greenstatus\quad\formalcheck

The map $T_0$ is continuous and increasing, fixes zero, and is strictly
below the identity on $(0,\infty)$ by \eqref{L7.6}.  Hence the sequence
$T_0^n(x^*)$ is decreasing and has a limit $\bar x\ge0$.  Continuity gives
$T_0(\bar x)=\bar x$; strict descent rules out $\bar x>0$, so
\begin{equation}
 T_0^n(x^*)\longrightarrow0.
 \eqL{7.7}\label{L7.7}
\end{equation}
\vspace{0.4\baselineskip}\par\noindent
Let $\widehat x>0$ be the binding cutoff from Property 3.  Equation
\eqref{L7.7} supplies an $N$ with $T_0^N(x^*)<\widehat x$.  Monotonicity of
$T_0^N$ then gives $T_0^N(x)<\widehat x$ simultaneously for every
$x\in[0,x^*]$.
\end{greenbox}

\begin{greenbox}{Positive-probability low-income entrance block:
\leanname{exists_positive_lowOffset_uniformly_enters_bindingRegion},
\leanname{lowOffsetBlockTransition_le_constant},
\leanname{lowLaborNeighborhoodBlock_uniformly_enters_bindingRegion},
\leanname{lowIncomeBlock_positive_and_uniformEntrance}}
\greenstatus\quad\formalcheck

For fixed $N$, continuity of \eqref{D55} and finite composition imply that
$\delta\mapsto T_\delta^N(x^*)$ is continuous.  The strict inequality in
\eqref{L7.7} therefore persists for some $\varepsilon>0$:
\begin{equation}
 T_\varepsilon^N(x^*)<\widehat x.
 \eqL{7.8}\label{L7.8}
\end{equation}
\par\smallskip
Both $T_\delta$ and the policy are increasing in their arguments.  Induction
in \eqref{D56} thus shows that, whenever $x\le x^*$ and
$0\le\delta_i\le\varepsilon$ for $i<N$,
\[
 Z_N(x;\boldsymbol\delta)\le T_\varepsilon^N(x^*)<\widehat x.
\]
Finally, lower-support reachability \eqref{A4} and independence give the
canonical event
\(\{l_i\le l_{\min}+\varepsilon:0\le i<N\}\) strictly positive
probability.  Thus the uniform low-income block needed by Aiyagari is now a
proved consequence of the solved policy, the Euler equation, and primitive
labor reachability; it is not a recurrence hypothesis.
\end{greenbox}

\begin{greenbox}{Actual household block has a common reset:
\leanname{supportedLowerLaborBlock_pos},
\leanname{continuous_absorbingSolvedBlockTransition},
\leanname{exists_oneStepSplitting_absorbingSolvedBlockTransition},
\leanname{kernelIterateOscillationMixing_of_blockSplitting}}
\greenstatus\quad\formalcheck

Put the supported labor law on the countable product history space and iterate
the actual absorbing random map for $N+1$ dates.  The first $N$ coordinates
form the proved positive-probability lower block.  They send every initial
$x\in S_{x^*}$ below $\widehat x$, where $A^*=0$.  Consequently, for the
additional shock $l_N$,
\begin{equation}
 X_{N+1}(x;l_0,\ldots,l_N)
   =w(l_N-l_{\min})\qquad\text{for every }x\in S_{x^*},
 \eqL{7.10}\label{L7.10}
\end{equation}
so the concrete finite-history transition satisfies the common-reset
splitting structure.  Its joint continuity is proved recursively.
\end{greenbox}

\begin{greenbox}{Finite-product representation and derived household mixing:
\leanname{finiteIIDBlockTransition_succ_last},
\leanname{measurable_finiteIIDBlockTransition},
\leanname{continuous_finiteIIDBlockTransition},
\leanname{integral_finiteIIDBlockTransition_eq_operator_iterate},
\leanname{fellerMarkovOperator_finiteIIDBlock_eq_iterate},
\leanname{finiteSupportedLowerLaborBlock_pos},
\leanname{exists_oneStepSplitting_finite_absorbingSolvedBlockTransition},
\leanname{absorbingSolvedKernel_finiteStepMixing}}
\greenstatus\quad\formalcheck

Let $\pi_l^{\otimes m}$ be the finite product law and let
$\Phi_m(x;l_0,\ldots,l_{m-1})$ denote chronological iteration of the actual
one-period random map.  Induction on $m$, using finite-product Fubini at the
successor step, proves the exact operator identity
\begin{equation}
 \int f\bigl(\Phi_m(x;\boldsymbol l)\bigr)\,
       d\pi_l^{\otimes m}(\boldsymbol l)
   =P_\kappa^m f(x).
 \eqL{7.11}\label{L7.11}
\end{equation}
The first $N$ coordinates of the concrete finite product are restricted to
the positive-probability lower block, while the last coordinate is free.
Equation \eqref{L7.10} makes $\Phi_{N+1}$ a common reset on this event.
Applying \eqref{L7.5} to the block kernel and then substituting
\eqref{L7.11} yields \eqref{D57} for the actual one-period household kernel.
Thus finite-step mixing is a theorem, not a hypothesis of Property 7.
\end{greenbox}

\begin{amberbox}{Property 7: \leanname{property07_unique_stable_invariant_distribution}}
\amberstatus\quad\formalcheck

Property 6 constructs \(x^*\); the preceding compact Feller theorem then
constructs an invariant probability on \(S_{x^*}\).  Lower-support
reachability, the low-income reset block, and \eqref{L7.11} prove the
finite-step oscillation condition \eqref{D57} for the actual restricted
household kernel.  The generic mixing theorem therefore gives
\begin{equation}
 \mathcal T_\kappa(\mu^*)=\mu^*,\qquad
 \mu^*\text{ is unique},\qquad
 \mathcal T_\kappa^n(\mu_0)\longrightarrow\mu^*\quad\forall\mu_0. \eqP{7.1}\label{P7.1}
\end{equation}

\missing{The stationarity and mixing argument itself is now foundationally
proved.  The numbered theorem remains amber at the paper level because it
still takes positive minimum-state consumption, local dominated
differentiation \eqref{A13}, and the Property 6 high-resource closure
\eqref{A14} as explicit analytic inputs.  Deriving those inputs from the
primitive utility and solved Bellman problem would upgrade the complete claim.}
\end{amberbox}

\begin{greenbox}{Continuity of unique invariant laws:
\leanname{continuous_integral_bcf_probabilityMeasure},
\leanname{continuous_fixedPoint_of_compact_unique},
\leanname{continuous_parameterizedMarkovLawEvolution_of_continuous_fellerOperator},
\leanname{continuous_invariantProbability_of_jointlyContinuous_evolution}}
\greenstatus\quad\formalcheck

Let $\Theta$ be a sequential parameter space, let $\mathcal P(S)$ be the
compact Hausdorff space of probability measures on a compact metric state
space, and let
\(\mathcal T:\Theta\times\mathcal P(S)\to\mathcal P(S)\) be jointly
continuous.  Suppose that for every $\theta$ there is a unique fixed point
$\mu_\theta$:
\begin{equation}
 \mathcal T(\theta,\mu_\theta)=\mu_\theta,
 \qquad
 \mathcal T(\theta,\nu)=\nu\Longrightarrow\nu=\mu_\theta.
 \eqL{7.12}\label{L7.12}
\end{equation}
Then $\theta\mapsto\mu_\theta$ is continuous.  Indeed, take
$\theta_n\to\theta$.  Compactness makes every cluster subsequence of
$\mu_{\theta_n}$ converge further to some $\bar\mu$.  Joint continuity and
\eqref{L7.12} imply
$\mathcal T(\theta,\bar\mu)=\bar\mu$; uniqueness gives
$\bar\mu=\mu_\theta$.  Thus every cluster point is $\mu_\theta$, and
compactness forces the whole sequence to converge.  Lean proves this first
for arbitrary compact fixed-point spaces and then specializes it to Markov
law evolution.  No continuity of the invariant-law selection is assumed.

The required joint continuity of law evolution is itself reduced to a
sup-norm statement about Feller operators.  For bounded continuous $f,g$ and
probability laws $\mu,\nu$,
\begin{equation}
 \left|\int g\,d\nu-\int f\,d\mu\right|
 \le \lVert g-f\rVert_\infty+
      \left|\int f\,d\nu-\int f\,d\mu\right|. \eqL{7.13}\label{L7.13}
\end{equation}
Thus joint integration is continuous in the sup-norm integrand and weak law.
Consequently, if $\theta\mapsto P_\theta f$ is sup-norm continuous for every
bounded continuous $f$, then
$(\theta,\mu)\mapsto\mathcal T_\theta(\mu)$ is jointly continuous.  The only
remaining model-specific continuity step is therefore continuity of the
solved household Feller operator, not an invariant-measure assumption.

\textbf{Remaining household application.}  To obtain continuity in
$(b,w,r)$, construct a common compact absorbing interval on each admissible
compact parameter set and prove joint continuity of the solved household law
evolution there.  The present primitives bundle prices into a single model,
so that parameterized-policy theorem is not yet available.
\end{greenbox}

For \(n\) states with fixed weights \(\omega_i\), define
\begin{equation}
 E_a(r)=\sum_{i=1}^n\omega_i\bigl(A_i(r)-\phi(r)\bigr). \eqP{8.1}\label{P8.1}
\end{equation}

\begin{greenbox}{Lemma 8.1: \leanname{finiteStateMeanAssets_continuous}}
\greenstatus\quad\formalcheck

If \(\phi\) and every \(A_i\) are continuous, each summand in (P8.1) is
continuous.  A finite sum of continuous functions is continuous.
\end{greenbox}

\begin{redbox}{Witness 8.2: \leanname{continuous_nonmonotone_finiteState_supply}}
\redstatus\quad\formalcheck

Take one state, \(\omega_1=1\), \(\phi=0\), and \(A_1(r)=r^2\).  Then
\(E_a(r)=r^2\), which is continuous, decreases from \(-1\) to \(0\), and
increases from \(0\) to \(1\); hence it is neither monotone nor antitone.

\missing{Show that this policy is optimal in an admissible finite-state
Aiyagari household model, or replace it with a certified CRRA witness.}
\end{redbox}

\begin{redbox}{Property 8: \leanname{property08_continuity_and_possible_nonmonotonicity}}
\redstatus\quad\formalcheck

Continuity of (P8.1) follows from Lemma 8.1, and the existential nonmonotone
component is supplied by Witness 8.2.

\missing{The continuity portion is rigorous, but the nonmonotonicity witness
must be an optimal policy and invariant distribution from an admissible
Aiyagari model.}
\end{redbox}

\begin{greenbox}{Lemma 9.1: \leanname{reciprocal_gap_diverges}}
\greenstatus\quad\formalcheck

As \(r\uparrow\lambda\), \(\lambda-r\downarrow0\) through positive values.
Continuity of subtraction and inversion therefore gives
\begin{equation}
 (\lambda-r)^{-1}\longrightarrow+\infty. \eqL{9.1}\label{L9.1}
\end{equation}
\end{greenbox}

\begin{amberbox}{Property 9: \leanname{property09_explosion_at_time_preference_rate}}
\amberstatus\quad\formalcheck

If eventually, as \(r\uparrow\lambda\),
\begin{equation}
 (\lambda-r)^{-1}\le E_a(r), \eqP{9.1}\label{P9.1}
\end{equation}
then Lemma 9.1 and order comparison of limits imply
\(E_a(r)\to+\infty\), namely (D35).

\missing{Derive (P9.1) from optimal saving, discounting, bounded shocks, and
risk-aversion conditions.}
\end{amberbox}

\begin{amberbox}{Property 10: \leanname{property10_uncertainty_raises_assets_near_benchmark}}
\amberstatus\quad\formalcheck

Let risky and certain asset supply be continuous at \(\lambda\), with
\(E_a^c(\lambda)<E_a^r(\lambda)\).  Strict inequality is open, so it persists
on some interval \((l,h)\) containing \(\lambda\).  Choosing
\(\epsilon=\lambda-l>0\) yields (D34).

\missing{Derive the strict gap at \(\lambda\) from a mean-preserving spread
and the associated invariant distributions.}
\end{amberbox}

\subsection{Properties 11--17: general equilibrium}

\begin{amberbox}{Property 11: \leanname{property11_factor_price_and_capital_signs}}
\amberstatus\quad\formalcheck

Let \(X=K-E_a^{IM}\).  Clearing at \(r^{IM}\) gives \(X(r^{IM})=0\), while
excess supply at \(r^{FI}\) gives \(X(r^{FI})<0\).  If \(X\) is strictly
decreasing, these facts force \(r^{IM}<r^{FI}\).  Equation (E7) gives
\(r^{FI}=P.\lambda\).  Since capital demand is strictly decreasing by (A11),
\begin{equation}
 r^{IM}<r^{FI}=P.\lambda,\qquad k^{FI}<k^{IM}. \eqP{11.1}\label{P11.1}
\end{equation}

\missing{Derive excess supply and the single-crossing condition from the
household and invariant-law results rather than assume them.}
\end{amberbox}

\begin{greenbox}{Lemma 12.1: \leanname{netSaving_eq_zero_of_goods_clear}}
\greenstatus\quad\formalcheck

From goods clearing \(C+\delta k=F(k)\), subtraction yields
\(F(k)-C-\delta k=0\).
\end{greenbox}

\begin{amberbox}{Property 12: \leanname{property12_saving_rate_sign}}
\amberstatus\quad\formalcheck

If \(P.s\) is strictly increasing, the capital ranking in (P11.1) gives
\(P.s(k^{FI})<P.s(k^{IM})\).  Lemma 12.1 applied to (E6) and (E7) makes net
saving zero in both equilibria.

\missing{Prove the hypotheses of Property 11 and the monotonicity of the
saving-rate formula from the production specification.}
\end{amberbox}

\begin{redbox}{Property 13: \leanname{property13_equilibrium_need_not_be_unique}}
\redstatus\quad\formalcheck

Use one state with weight one, \(\phi(r)=1\), \(A(r)=r^2\), and \(K(r)=0\).
Then (P8.1) is \(E_a(r)=r^2-1\), which is nonmonotone, and both \(r=-1\) and
\(r=1\) satisfy \(K(r)=E_a(r)=0\).

\missing{Certify the policy as optimal, the weights as invariant, and the
capital-demand curve as arising from admissible production.  The current
witness is not an Aiyagari equilibrium.}
\end{redbox}

\begin{greenbox}{Lemma 14.1: \leanname{marketClearing_exists_of_sign_change}}
\greenstatus\quad\formalcheck

If \(X=K-E_a\) is continuous on \([r_L,r_H]\), with
\(X(r_L)\ge0\ge X(r_H)\), the intermediate value theorem supplies
\(r\in[r_L,r_H]\) with \(X(r)=0\), equivalently (D31).
\end{greenbox}

\begin{redbox}{Property 14: \leanname{property14_existence_under_alternative_debt_limits}}
\redstatus\quad\formalcheck

Apply Lemma 14.1 to \((K_{PV},E_{PV})\) on a positive interval and to
\((K_f,E_f)\) on a negative interval.  The endpoint restrictions place the
first root above zero and the second below zero.

\missing{Define these curves from the present-value and fixed-debt household
models and prove the endpoint signs.  As stated they are arbitrary curves, so
the theorem does not establish existence under the named debt regimes.}
\end{redbox}

\begin{greenbox}{Lemma 15.1: \leanname{assetSupply_eq_capitalDemand_at_equilibrium}}
\greenstatus\quad\formalcheck

Market clearing is \(K(r)=E_a(r)\) by (D31); symmetry gives
\(E_a(r)=K(r)\).
\end{greenbox}

\begin{greenbox}{Property 15: \leanname{property15_general_equilibrium_disciplines_saving}}
\greenstatus\quad\formalcheck

If \(K(r_p)<E_a(r_p)\), then equality cannot hold at \(r_p\), so that partial
return is not an equilibrium.  At a clearing return \(r_e\), Lemma 15.1 gives
\(E_a(r_e)=K(r_e)\).
\end{greenbox}

\begin{amberbox}{Property 16: \leanname{property16_more_borrowing_reduces_capital}}
\amberstatus\quad\formalcheck

Let $E_a^0$ and $E_a^1$ be the asset-supply schedules under the tight and
relaxed limits.  Lean now assumes only the economically relevant local
comparison $E_a^1(r_0)<E_a^0(r_0)$ at the old clearing return, rather than an
artificial constant translation of the entire schedule.  Since
$K(r_0)=E_a^0(r_0)$, new excess demand at $r_0$ is strictly positive; at the
new root it is zero.  If new excess demand is strictly decreasing, this forces
$r_0<r_1$.  Strictly decreasing capital demand then gives
$K(r_1)<K(r_0)$.

\missing{Derive $E_a^1(r_0)<E_a^0(r_0)$ from the two solved policies and their
respective invariant laws under a relaxed borrowing constraint.}
\end{amberbox}

\begin{amberbox}{Property 16b: \leanname{property16_more_borrowing_reduces_saving_rate}}
\amberstatus\quad\formalcheck

If \(r_0<r_1\), decreasing primitive capital demand gives
$K(r_1)<K(r_0)$; applying the actual technology saving-rate function
$s_P(k)=\delta k/f(k)$ preserves the inequality when $s_P$ is strictly
increasing.

\missing{Connect the assumed return comparison to the borrowing-limit model
and derive saving-rate monotonicity from production.}
\end{amberbox}

\begin{greenbox}{Lemma 17.1: \leanname{effectiveBorrowingLimit_eq_natural_of_slack}}
\greenstatus\quad\formalcheck

For \(r>0\), (D3) is the minimum of \(b\) and \(\phi^{nat}\).  Slackness
(D36) says \(\phi^{nat}\le b\), so the minimum is \(\phi^{nat}\).
\end{greenbox}

\begin{greenbox}{Solved-household equality lemmas:
\leanname{effectiveBorrowingLimit_admissible},
\leanname{solvedBoundedHousehold_behavior_eq_of_limit_eq},
\leanname{solvedNextExcess_eq_of_limit_eq},
\leanname{solvedResourceKernel_eq_of_limit_eq}}
\greenstatus\quad\formalcheck

For $r>0$, $\phi^{eff}\le\phi^{nat}$ implies
$z_{\min}=wl_{\min}-r\phi^{eff}\ge0$, so the effective limit is admissible.
If $\phi_1=\phi_2$, the two concrete Bellman operators are definitionally the
same.  Banach's construction, the unique optimizer, optimal consumption, and
the policy-induced transition therefore agree.  Applying the identical
random transition to the same labor law gives the same Markov kernel.  Proof
witnesses for admissibility lie in $\mathsf{Prop}$ and cannot change any of
these economic objects.
\end{greenbox}

\begin{amberbox}{Property 17: \leanname{property17_no_effect_when_fixed_limit_slack}}
\amberstatus\quad\formalcheck

If both \(b_1\) and \(b_2\) are slack, Lemma 17.1 makes both effective limits
equal \(\phi^{nat}\).  The theorem now applies the preceding concrete lemmas,
not an arbitrary behavior function, and proves
\begin{equation}
 V_{b_1}=V_{b_2},\quad A_{b_1}=A_{b_2},\quad c_{b_1}=c_{b_2},\quad
 T_{b_1}=T_{b_2},\quad P_{b_1}=P_{b_2}. \eqP{17.1}\label{P17.1}
\end{equation}
Consequently a probability measure is invariant for $P_{b_1}$ if and only if
it is invariant for $P_{b_2}$.  This establishes the fixed-price household
and stationary-law content of the slackness claim.

\missing{Upgrade the result to an explicit bijection between complete
stationary recursive equilibrium records with statutory limits $b_1$ and
$b_2$, including the bookkeeping change in the model's stored $b$ field and
unchanged aggregate market-clearing equations.}
\end{amberbox}

\subsection{Properties 18--20: growth, debt, and money}

Define exact growth and normalization by
\begin{equation}
 x_t=x_0(1+g)^t,\qquad \widetilde x_t=\frac{x_t}{(1+g)^t}. \eqP{18.1}\label{P18.1}
\end{equation}

\begin{greenbox}{Lemma 18.1: \leanname{normalizedPath_eq}}
\greenstatus\quad\formalcheck

If \(1+g\ne0\), substitution into (P18.1) cancels the nonzero factor and
gives \(\widetilde x_t=x_0\) at every date.
\end{greenbox}

\begin{amberbox}{Property 18: \leanname{property18_balanced_growth_extension}}
\amberstatus\quad\formalcheck

Assume \(\beta>0\), \(1+g\ne0\), exact growth for earnings, resources, assets,
and consumption, and
\begin{equation}
 \beta(1+r^{IM})<(1+g)^\gamma. \eqP{18.2}\label{P18.2}
\end{equation}
Since \(1+\lambda(\beta)=1/\beta\), division of (P18.2) by \(\beta>0\) gives
\(r^{IM}<r_g^{FI}\) from (D37).  Lemma 18.1 makes all four normalized paths
constant at their initial levels.

\missing{Derive the Euler wedge and balanced-growth paths from the growing
economy rather than assume them.}
\end{amberbox}

\begin{greenbox}{Lemma 19.1: \leanname{debt_budget_translation}}
\greenstatus\quad\formalcheck

For gross assets \(a^{net}+d\), expand (D38):
\begin{equation}
 R_d(w,l,r,d,a^{net}+d)-d=wl+(1+r)a^{net}. \eqL{19.1}\label{L19.1}
\end{equation}
The \(-rd+(1+r)d-d\) terms cancel.
\end{greenbox}

\begin{greenbox}{Lemma 19.2: \leanname{debt_budget_iff_net_budget}}
\greenstatus\quad\formalcheck

Adding \(d\) to next-period net assets and using (L19.1) shows
\[
 c+(a'_{net}+d)=R_d(w,l,r,d,a_{net}+d)
 \Longleftrightarrow c+a'_{net}=wl+(1+r)a_{net}.
\]
\end{greenbox}

\begin{greenbox}{Lemma 19.3: \leanname{taxAdjusted_floor_iff_net_natural_floor}}
\greenstatus\quad\formalcheck

For \(r>0\), definitions (D2) and (D40) give
\(\underline a_d-d=-\phi^{nat}\).  Consequently
\begin{equation}
 \underline a_d\le a^{net}+d\Longleftrightarrow
 -\phi^{nat}\le a^{net}. \eqL{19.2}\label{L19.2}
\end{equation}
\end{greenbox}

Let gross and net paths be \(a_t=a_t^{net}+d\) and
\(a_t^{net}=a_t-d\).  Debt-adjusted feasibility combines the gross budget
with \(a_t\ge\underline a_d\); debt-free feasibility combines the net budget
with \(a_t^{net}\ge-\phi^{nat}\).

\begin{greenbox}{Lemma 19.4: \leanname{debtAdjustedFeasible_iff_net}}
\greenstatus\quad\formalcheck

Apply Lemma 19.2 at every date to the budget equations and Lemma 19.3 at every
date to the floors.  Both implications follow componentwise, giving an
equivalence of the two feasible-path predicates.
\end{greenbox}

\begin{greenbox}{Debt Bellman equivalence:
\leanname{debtAdjustedBellmanObjective_grossAssetOfShifted},
\leanname{debtAdjustedBellmanOperator_eq_bellmanOperator},
\leanname{debtAdjustedBellmanFixedPoint_of_fixedPoint},
\leanname{debtAdjustedPolicy_optimal_of_optimal},
\leanname{debtAdjustedResourceTransition_policy}}
\greenstatus\quad\formalcheck

Let \(g=g_d(\widehat a)=\widehat a-\phi+d\).  Equations (E10)--(E12) and
(D38) give
\begin{align}
 c_d(z,g_d(\widehat a))&=z-\widehat a, \notag\\
 T_d(g_d(\widehat a),l')&=T_{M,\phi}(\widehat a,l'), \notag\\
 Q_d(V;z,g_d(\widehat a))&=Q_{M,\phi}(V;z,\widehat a). \eqL{19.3}\label{L19.3}
\end{align}
Moreover \(g_d\) maps \([0,z]\) bijectively onto
\([d-\phi,z+d-\phi]\), with inverse \(g\mapsto g+\phi-d\).  Hence the two
sets of objective values are equal, so their suprema are equal:
\begin{equation}
 \mathcal B_dV=\mathcal B_{M,\phi}V. \eqL{19.4}\label{L19.4}
\end{equation}
Thus every debt-free Bellman fixed point is a gross-debt fixed point, and the
image under \(g_d\) of every debt-free maximizing policy maximizes the actual
gross-debt objective.  This proves household optimality rather than assuming
it as a neutrality premise.
\end{greenbox}

\begin{greenbox}{Complete-equilibrium equivalence:
\leanname{taxAdjustedToDebtFree_debtFreeToTaxAdjusted},
\leanname{debtFreeToTaxAdjusted_taxAdjustedToDebtFree},
\leanname{taxAdjustedAssetFloor_eq_debt_sub_natural}}
\greenstatus\quad\formalcheck

From a debt-free record satisfying (E13), retain \((V,A,\mu)\) and use the
gross policy \(g_d\circ A\).  Bellman optimality follows from (L19.3)--(L19.4),
and the transition equality in (L19.3) preserves the same invariant law.
Integrability is preserved by addition of the constant \(d\), while
probability of \(\mu\) and (E13) give
\begin{equation}
 \int g_d(A(z))\,d\mu
 =\int(A(z)-\phi)\,d\mu+d\mu(\mathbb R)=0+d=d. \eqL{19.5}\label{L19.5}
\end{equation}
This constructs a complete record satisfying (E14).  Subtracting \(d\)
constructs the reverse map; (L19.5) then yields zero net-asset supply.  The two
maps preserve the household and measure fields exactly and are inverse (proof
fields are propositions and hence proof-irrelevant).  Finally, for \(r>0\),
direct algebra gives
\(\underline a_d=d-\phi^{nat}\), so when \(\phi=\phi^{nat}\) the gross
feasible interval in (E10) has precisely the tax-adjusted natural floor.
\end{greenbox}

\begin{greenbox}{Witness 19.6: \leanname{fixed_floor_not_debt_neutral_witness}}
\greenstatus\quad\formalcheck

Choose \(b=0,d_0=0,d_1=1,a^{net}=-1/2\).  Then
\(-b=0\le a^{net}+d_1=1/2\), whereas
\(0\nleq a^{net}+d_0=-1/2\).  Thus a fixed gross floor need not be preserved
by debt translation.
\end{greenbox}

\begin{greenbox}{Property 19: \leanname{property19_debt_neutrality_depends_on_constraint}}
\greenstatus\quad\formalcheck

Assume the household primitives satisfy (A1), \(r>0\), and
\(\phi=\phi^{nat}\).  The explicit map from (E13) to
(E14) is bijective because the two preceding inverse identities hold.  It
preserves Bellman values, optimal policies, consumption, resource transitions,
invariant distributions, and asset-market clearing, while (L19.2) preserves
every feasible path.  Direct cancellation gives
\((a_t^{net}+d)-d=a_t^{net}\).  Witness 19.6 shows why the conclusion fails
for a fixed gross borrowing floor that is not adjusted with government debt.
Thus the tax-adjusted constraint gives neutrality of the complete stationary
equilibrium set, whereas a fixed constraint need not.
\end{greenbox}

\begin{greenbox}{Lemma 20.1: \leanname{sharpPrice_pos}}
\greenstatus\quad\formalcheck

If \(r^\#,w,l_{\min}>0\), then (D42) is a positive number divided by a
positive number, hence \(p^\#>0\).
\end{greenbox}

\begin{greenbox}{Lemma 20.2: \leanname{one_div_sharpPrice_eq_naturalLimit}}
\greenstatus\quad\formalcheck

When all three factors are nonzero,
\begin{equation}
 \frac1{p^\#}=\frac{wl_{\min}}{r^\#}=\phi^{nat}(w,l_{\min},r^\#). \eqL{20.1}\label{L20.1}
\end{equation}
\end{greenbox}

\begin{greenbox}{Lemma 20.3: \leanname{monetary_constraint_slack_above_sharp_return}}
\greenstatus\quad\formalcheck

Suppose \(0<r^\#<r\), \(w,l_{\min},p>0\), and \(p<p^\#\).  Since
\(wl_{\min}/r\) decreases in positive \(r\),
\[
 \phi^{nat}(r)<\phi^{nat}(r^\#)=1/p^\#<1/p.
\]
Thus the monetary fixed limit \(1/p\) is slack, and Lemma 17.1 makes the
effective limit equal the natural limit.
\end{greenbox}

\begin{greenbox}{Stationary mean equality:
\leanname{SolvedStationaryHousehold.meanAssets_eq_of_limit_eq_unique}}
\greenstatus\quad\formalcheck

Let \(S_1,S_2\) be the concrete records (E15) for the same household model.
If \(\phi_1=\phi_2\), the solved policy and full-state kernel coincide by the
equal-limit household and kernel lemmas.  Hence \(S_1.\mu\) is invariant for
the second kernel.  If the second invariant probability is unique, then
\(S_1.\mu=S_2.\mu\).  Equality of both the integrand and measure in (E15)
therefore proves
\begin{equation}
 \mathfrak A(S_1)=\mathfrak A(S_2). \eqL{20.2}\label{L20.2}
\end{equation}
\end{greenbox}

\begin{greenbox}{Lemma 20.4: \leanname{monetaryMeanAssets_eq_natural_of_slack}}
\greenstatus\quad\formalcheck

Let \(S_m\) be the actual monetary stationary record (E16) and \(S_n\) a
natural-limit stationary record for the same model.  If the monetary limit is
slack and \(S_n\)'s invariant probability is unique, then the two effective
limits agree.  The preceding theorem—not substitution into an arbitrary
schedule—gives
\begin{equation}
 \mathfrak A(S_m)=\mathfrak A(S_n). \eqL{20.3}\label{L20.3}
\end{equation}
\end{greenbox}

\begin{amberbox}{Property 20: \leanname{property20_monetary_return_upper_bound}}
\amberstatus\quad\formalcheck

Let \(S_m\) be a solved monetary equilibrium (E16) for model \(M\), let
\(S_n\) be its natural-limit stationary comparator, and let \(S^\#\) be a
natural-limit zero-asset equilibrium (E17) for a model \(M^\#\) with the same
household environment except for the return.  Assume \(S_n\)'s invariant law
is unique and the two concrete policy integrals obey the pairwise return
ordering
\begin{equation}
 M^\#.r<M.r\quad\Longrightarrow\quad
 \mathfrak A(S^\#)<\mathfrak A(S_n). \eqP{20.1}\label{P20.1}
\end{equation}
If \(M.r>M^\#.r\), Lemma 20.3 makes the monetary limit slack and Lemma 20.4
gives \(\mathfrak A(S_m)=\mathfrak A(S_n)\).  Equations (E16)--(E17) make both
sides of (P20.1) zero, an impossibility.  Hence \(M.r\le M^\#.r\).

\missing{Derive the full-state uniqueness premise from Property 7's compact
invariant law and derive the pairwise ordering (P20.1) from the natural-limit
asset-supply results.  The former arbitrary supply schedule and representation
hypothesis have been eliminated.}
\end{amberbox}

\section{Declaration coverage and adequacy summary}

Every declaration below is kernel-checked and independent of
\texttt{sorryAx}.  ``Upgrade step'' is therefore about economic or analytic
adequacy, never logical validity.

{\scriptsize
\setlength{\tabcolsep}{2.5pt}
\renewcommand{\arraystretch}{1.08}
\begin{longtable}{>{\raggedright\arraybackslash}p{0.32\textwidth}
                  >{\raggedright\arraybackslash}p{0.08\textwidth}
                  >{\raggedright\arraybackslash}p{0.09\textwidth}
                  >{\raggedright\arraybackslash}p{0.22\textwidth}
                  >{\raggedright\arraybackslash}p{0.20\textwidth}}
\toprule
Lean declaration & Soundness & Tier & Reason & Upgrade step \\
\midrule
\endfirsthead
\toprule Lean declaration & Soundness & Tier & Reason & Upgrade step \\
\midrule
\endhead
\leanname{consumption_nonneg} & Checked & Green & Directly unfolds feasible policy and consumption. & None. \\
\leanname{market_clears} & Checked & Green & Directly unfolds equilibrium clearing. & None. \\
\leanname{minimumIncome_shifted_budget} & Checked & Green & Exact budget-shift algebra. & None. \\
\leanname{discountedShiftedAssets_antitone} & Checked & Green & Feasibility signs the discounted recursion. & None. \\
\leanname{tendsto_naturalDebtLimit_div_pow} & Checked & Green & Geometric-limit argument under \(r>0\). & None. \\
\leanname{property01_naturalDebt_iff_noPonzi} & Checked & Green & Canonical i.i.d. proof uses predictability, independence, and positive-probability low-income blocks. & None. \\
\leanname{property01_naturalDebt_iff_noPonzi_ae} & Checked & Green & Exact almost-sure alias of the canonical theorem. & None. \\
\leanname{natural_bound_does_not_imply_noPonzi_without_feasibility} & Checked & Green & Explicit deterministic counterexample. & None. \\
\leanname{noPonzi_does_not_imply_natural_bound_without_feasibility} & Checked & Green & Explicit deterministic counterexample. & None. \\
\leanname{beta_mul_one_add_timePreferenceRate} & Checked & Green & Exact identity from the time-preference definition. & None. \\
\leanname{canonicalLabor_map_eq} & Checked & Green & A product-measure coordinate has the primitive labor law. & None. \\
\leanname{measure_lowerLaborBlock} & Checked & Green & Finite cylinder probability is the product of marginal probabilities. & None. \\
\leanname{lowerLaborBlock_pos} & Checked & Green & Reachability makes every factor, hence the finite product, positive. & None. \\
\leanname{exists_compactMaximum_eq} & Checked & Green & Extreme-value theorem on a nonempty compact action space. & None. \\
\leanname{compactMaximum_le} & Checked & Green & Every value is below the attained supremum. & None. \\
\leanname{abs_compactMaximum_sub_le} & Checked & Green & Attained compact maxima obey the uniform perturbation bound. & None. \\
\leanname{compactMaximum_lipschitz} & Checked & Green & The perturbation estimate is exactly one-Lipschitz continuity. & None. \\
\leanname{continuous_compactMaximum} & Checked & Green & Lipschitz maps are continuous. & None. \\
\leanname{continuousFamilyArgmax_isMax}, \leanname{continuous_continuousFamilyArgmax} & Checked & Green & Compactness, joint objective continuity, and optimizer uniqueness prove Berge continuity by a subsequence argument. & Specialize uniqueness to each solved household objective. \\
\leanname{maximizeObjective_norm_sub_le} & Checked & Green & Statewise compact maximization is sup-norm nonexpansive. & None. \\
\leanname{abs_integral_le_uniform_bound} & Checked & Green & Probability normalization bounds an integral by its uniform envelope. & None. \\
\leanname{abs_integral_bcf_sub_le} & Checked & Green & Probability integration is sup-norm nonexpansive. & None. \\
\leanname{expectedContinuationBCF_norm_sub_le} & Checked & Green & Continuous Markov continuation preserves bounded continuity and is nonexpansive. & None. \\
\leanname{continuous_boundedContinuousFamily}, \leanname{continuous_compactParameterizedExpectation_fixedValue}, \leanname{continuous_compactParameterizedExpectation} & Checked & Green & Heine--Cantor and probability nonexpansiveness give sup-norm continuity on a common compact state space. & Establish a common state for the global household family. \\
\leanname{continuousBellmanOperator_norm_sub_le} & Checked & Green & Reward cancels and discounting gives the beta bound. & None. \\
\leanname{continuousBellmanOperator_contracting} & Checked & Green & The preceding estimate and beta below one prove contraction. & None. \\
\leanname{parameterizedContractionFixedPoint_isFixedPt}, \leanname{dist_parameterizedContractionFixedPoint_le}, \leanname{continuous_parameterizedContractionFixedPoint} & Checked & Green & The residual estimate \eqref{L2.5a} proves continuity of fixed points for jointly continuous uniformly contracting families. & Prove joint continuity of the concrete \((b,w,r)\)-Bellman family. \\
\leanname{continuous_continuousBellmanOperator_reward_value}, \leanname{continuous_boundedBellmanValue_reward} & Checked & Green & The concrete bounded Bellman map is jointly continuous in reward and candidate value, so its fixed point is continuous in reward. & Extend the operator calculation to price-dependent rewards and transitions. \\
\leanname{continuous_compactParameterizedBellmanOperator}, \leanname{continuous_compactParameterizedBellmanValue}, \leanname{continuous_compactParameterizedSolvedObjective}, \leanname{continuous_compactParameterizedSolvedObjective_apply}, \leanname{continuous_compactParameterizedSolvedPolicy} & Checked & Green & On common compact spaces, primitive reward/transition continuity implies solved value, objective, and unique-policy continuity. & Prove global truncation equivalence or use a weighted space. \\
\leanname{boundedBellmanValue_isFixedPoint} & Checked & Green & Generic Banach fixed-point construction. & None. \\
\leanname{boundedBellmanValue_unique} & Checked & Green & Generic Banach uniqueness. & None. \\
\leanname{shareConsumption_nonneg} & Checked & Green & Nonnegative resources and saving share imply nonnegative consumption. & None. \\
\leanname{continuous_shareNextExcess} & Checked & Green & Arithmetic and nonnegative clipping preserve continuity. & None. \\
\leanname{boundedAiyagariValue_isFixedPoint} & Checked & Green & Banach constructs the bounded continuous fixed point. & None. \\
\leanname{boundedAiyagariValue_unique} & Checked & Green & Banach fixed-point uniqueness. & None. \\
\leanname{boundedAiyagariPolicy_isMax} & Checked & Green & Compactness and continuity produce a pointwise maximizer. & None. \\
\leanname{boundedAiyagariValue_concave} & Checked & Green & Bellman iteration and closedness of concavity pass concavity to the fixed point. & None. \\
\leanname{boundedOptimalAsset_unique} & Checked & Green & Primitive strict utility concavity and continuation concavity give a unique optimizer. & None. \\
\leanname{continuous_boundedAiyagariAssetPolicy} & Checked & Green & Compact-subsequence maximum-theorem proof with optimizer uniqueness. & None. \\
\leanname{hasDerivAt_of_twoSided_slope_sandwich} & Checked & Green & Two-sided optimality slope bounds imply the envelope derivative by squeezing. & Apply it to the solved Bellman objective. \\
\leanname{hasDerivAt_of_concave_local_lowerSupport} & Checked & Green & A touching differentiable lower support pins down both one-sided derivatives of a concave value. & None. \\
\leanname{LocalDominatedDerivative.hasDerivAt_integral} & Checked & Green & The expectation derivative follows from the explicit domination data in \eqref{A13}. & None. \\
\leanname{HasDerivAt.nonpos_of_eventually_right_max} & Checked & Green & Right-feasible optimality makes the derivative nonpositive. & None. \\
\leanname{boundedAiyagariValue_eq_optimalAssetObjective} & Checked & Green & The fixed point equals its objective at the constructed optimizer. & None. \\
\leanname{boundedAiyagari_envelope_of_consumption_pos} & Checked & Green & Concavity and a feasible fixed-action support prove \eqref{P2.1}. & None. \\
\leanname{boundedAssetContinuation_hasDerivAt} & Checked & Green & Dominated differentiation derives the continuation derivative. & None. \\
\leanname{boundedAssetObjective_hasDerivAt_at_policy} & Checked & Green & Chain, sum, and parametric-integral rules derive the objective derivative. & None. \\
\leanname{boundedAiyagari_euler_inequality} & Checked & Green & Boundary and interior first-order conditions prove \eqref{P2.2}. & None. \\
\leanname{property02_value_and_policy_regularity} & Checked & Amber & The constructed solution now includes the envelope and conditional Euler system in addition to existence, uniqueness, concavity, optimality, and state continuity. & Establish concrete Bellman/operator and optimizer parameter continuity, global domination, and weighted Inada theory. \\
\leanname{boundedAiyagariAssetPolicy_zero_at_minimum} & Checked & Green & The Euler contradiction and impatience prove zero saving at the minimum state. & None, conditional on the explicitly stated analytic closure conditions. \\
\leanname{boundedAiyagari_marginal_le_minimum_of_policy_pos} & Checked & Green & Interior Euler equality, policy monotonicity, and decreasing marginal utility prove \eqref{L3.1}. & None, conditional on the explicitly stated analytic closure conditions. \\
\leanname{boundedAiyagari_binding_cutoff} & Checked & Green & Continuity and the strict impatience gap turn the marginal bound into a positive binding interval. & None, conditional on the explicitly stated analytic closure conditions. \\
\leanname{property03_binding_region_at_low_resources} & Checked & Amber & A nondegenerate binding interval is now derived from the solved policy and Euler system. & Derive positive minimum consumption and global dominated differentiation from primitives. \\
\leanname{property03_inada_exception} & Checked & Red & Positivity is assumed; the Inada condition is unused. & Derive interiority from Inada utility, feasibility, and optimality. \\
\leanname{monotone_boundedAiyagariAssetPolicy} & Checked & Green & Revealed preference and strict utility concavity prove asset monotonicity. & None. \\
\leanname{monotone_boundedAiyagariConsumption} & Checked & Green & Strict optimality and continuation concavity prove consumption monotonicity. & None. \\
\leanname{boundedAiyagariAssetPolicy_increment_bounds} & Checked & Green & The two monotonicity results imply global finite-difference bounds between zero and one. & None. \\
\leanname{property04_monotonicity_and_one_for_one_response} & Checked & Amber & Concrete weak monotonicity, continuity, and weak slope bounds are proved. & Derive strict derivatives and combine with the nondegenerate cutoff. \\
\leanname{consumptionShiftedDown_of_policyShiftedUp} & Checked & Green & Pointwise subtraction reverses the policy shift. & None. \\
\leanname{integral_mono_of_convexOrder_concave} & Checked & Green & Convex order reverses expectation order for integrable concave tests. & None. \\
\leanname{continuationSavingGap_mono_of_meanPreservingSpread} & Checked & Green & Miller-class shifted differences and convex order prove \eqref{L5.2}. & None. \\
\leanname{optimalConsumption_antitone_of_meanPreservingSpread_sameContinuation} & Checked & Green & Optimality, uniqueness, and \eqref{L5.2} prove the one-step policy comparison. & None. \\
\leanname{property05_mean_preserving_spread} & Checked & Amber & Genuine convex order now yields lower consumption and higher saving for a common Miller-class continuation value. & Prove Bellman preservation and monotonicity of the difference between the two law-specific fixed points. \\
\leanname{EventualRelativeRiskAversionBound.marginal_ratio_le_rpow} & Checked & Green & Integrating the primitive differential RRA bound proves \eqref{L6.1}. & None. \\
\leanname{EventualRelativeRiskAversionBound.tendsto_marginal_ratio_add_const} & Checked & Green & Marginal monotonicity, \eqref{L6.1}, and squeezing prove \eqref{L6.2}. & None. \\
\leanname{rra_marginal_ratio_bounded_gap_tendsto_one} & Checked & Green & The bounded labor gap and \eqref{L6.2} prove \eqref{L6.3}. & None. \\
\leanname{bounded_resource_dynamics_of_policy_tendsto_atTop} & Checked & Green & Euler, envelope, RRA, and value concavity prove drift in the unbounded-policy branch. & None, conditional on labeled closures \eqref{A13}--\eqref{A14}. \\
\leanname{property06_bounded_resource_dynamics} & Checked & Amber & The solved-policy drift theorem is proved by the bounded/unbounded policy dichotomy. & Derive high-resource consumption divergence \eqref{A14} from primitives. \\
\leanname{feller_solvedResourceTransition} & Checked & Green & Joint continuity and compact-support parametric integration prove Feller continuity. & None. \\
\leanname{monotone_solvedResourceTransition} & Checked & Green & Policy monotonicity and positive gross return preserve the state order. & None. \\
\leanname{exists_invariantProbability_of_feller_compact} & Checked & Green & Concrete Cesàro averages, compactness, Feller continuity, and the telescoping residual prove invariant-law existence. & None. \\
\leanname{solvedHousehold_invariantProbability_exists} & Checked & Green & Property 6 restricts the solved Feller household to a compact absorbing interval, where Krylov--Bogoliubov applies. & None, conditional on the labeled Property 6 closures. \\
\leanname{exists_unique_invariant_and_tendsto_of_feller_mixing} & Checked & Green & Oscillation contraction yields geometric test-function convergence and uniqueness. & None, conditional on the explicit local kernel condition \eqref{D53}. \\
\leanname{tendsto_bcfOscillation_iterate_of_finiteStepMixing} & Checked & Green & Complete contracting blocks dominate the finitely many remainder steps, proving \eqref{L7.9}. & None. \\
\leanname{exists_unique_invariant_and_tendsto_of_feller_finiteStepMixing} & Checked & Green & Finite-step contraction yields uniqueness and convergence at every date of the original chain. & None, conditional on \eqref{D57}. \\
\leanname{kernelOscillationMixing_of_oneStepSplitting} & Checked & Green & A positive common-reset event proves \eqref{L7.5} with coefficient below one. & None. \\
\leanname{minimumLaborTransition_strictly_decreases} & Checked & Green & Euler equality, the envelope theorem, monotone consumption, and impatience prove \eqref{L7.6}. & None, conditional on the labeled dominated-differentiation input. \\
\leanname{tendsto_iterate_zero_of_monotone_strictDecrease} & Checked & Green & Monotone convergence and continuity identify zero as the only possible iterate limit. & None. \\
\leanname{minimumLaborBlock_uniformly_enters_bindingRegion} & Checked & Green & The minimum-income iterates converge to zero uniformly over every bounded initial interval by monotonicity. & None. \\
\leanname{exists_positive_lowOffset_uniformly_enters_bindingRegion} & Checked & Green & Continuity makes strict entrance persist on an open relative neighborhood of the lower labor endpoint. & None. \\
\leanname{lowOffsetBlockTransition_le_constant} & Checked & Green & Coordinatewise transition monotonicity bounds every low-offset history by the constant upper-offset path. & None. \\
\leanname{lowIncomeBlock_positive_and_uniformEntrance} & Checked & Green & Lower-support reachability and i.i.d. block probabilities combine with \eqref{L7.8}. & None. \\
\leanname{exists_oneStepSplitting_absorbingSolvedBlockTransition} & Checked & Green & The supported product block has positive probability and the following shock coalesces all initial states as in \eqref{L7.10}. & None. \\
\leanname{kernelIterateOscillationMixing_of_blockSplitting} & Checked & Green & A block random-map representation plus common reset proves finite-step mixing. & None, conditional on the displayed representation equality. \\
\leanname{integral_finiteIIDBlockTransition_eq_operator_iterate} & Checked & Green & Finite-product Fubini proves the exact block expectation identity \eqref{L7.11}. & None. \\
\leanname{fellerMarkovOperator_finiteIIDBlock_eq_iterate} & Checked & Green & Repackages \eqref{L7.11} as equality of the block Markov operator and the one-period operator iterate. & None. \\
\leanname{exists_oneStepSplitting_finite_absorbingSolvedBlockTransition} & Checked & Green & The concrete finite product contains a positive-probability common-reset block. & None. \\
\leanname{absorbingSolvedKernel_finiteStepMixing} & Checked & Green & The finite reset block and \eqref{L7.11} derive \eqref{D57} for the actual solved household kernel. & None, conditional on the displayed local analytic inputs. \\
\leanname{property07_unique_stable_invariant_distribution} & Checked & Amber & Compact-Feller existence, concrete household mixing, uniqueness, and convergence from every initial law are all proved. & Derive positive minimum-state consumption and closures \eqref{A13}--\eqref{A14} from primitives. \\
\leanname{continuous_integral_bcf_probabilityMeasure} & Checked & Green & The decomposition \eqref{L7.13} proves joint continuity in the sup-norm integrand and weak probability law. & None. \\
\leanname{continuous_fixedPoint_of_compact_unique} & Checked & Green & Compactness, joint continuity, and pointwise fixed-point uniqueness prove continuity of the fixed-point selection. & None. \\
\leanname{continuous_parameterizedMarkovLawEvolution_of_continuous_fellerOperator} & Checked & Green & Sup-norm continuity of each parameterized Feller operator implies joint continuity of law evolution. & Prove this operator continuity for the solved household. \\
\leanname{continuous_invariantProbability_of_jointlyContinuous_evolution} & Checked & Green & Specializes the compact fixed-point argument to parameterized Markov law evolution. & Prove joint law-evolution continuity for the solved household on common parameter-local absorbing intervals. \\
\leanname{finiteStateMeanAssets_continuous} & Checked & Green & A finite weighted sum preserves continuity. & None. \\
\leanname{continuous_nonmonotone_finiteState_supply} & Checked & Red & The \(r^2\) policy is not certified optimal. & Construct an admissible household-model witness. \\
\leanname{property08_continuity_and_possible_nonmonotonicity} & Checked & Red & It inherits the uncertified \(r^2\) witness. & Verify policy optimality, invariance, and admissible primitives. \\
\leanname{reciprocal_gap_diverges} & Checked & Green & Standard reciprocal limit. & None. \\
\leanname{property09_explosion_at_time_preference_rate} & Checked & Amber & Explosion uses a supplied divergent lower bound. & Derive the bound from the Euler equation and stationary law. \\
\leanname{property10_uncertainty_raises_assets_near_benchmark} & Checked & Amber & Continuity propagates an assumed precautionary gap. & Prove the gap from a mean-preserving spread. \\
\leanname{property11_factor_price_and_capital_signs} & Checked & Amber & Excess-demand signs and single crossing are inputs. & Derive them from household supply and production. \\
\leanname{netSaving_eq_zero_of_goods_clear} & Checked & Green & Goods clearing yields the identity algebraically. & None. \\
\leanname{property12_saving_rate_sign} & Checked & Amber & Saving comparisons use supplied monotonicity. & Derive factor-price and saving-rate shapes from primitives. \\
\leanname{property13_equilibrium_need_not_be_unique} & Checked & Red & Arbitrary aggregation and zero capital demand are used. & Certify two market-clearing equilibria of an admissible model. \\
\leanname{marketClearing_exists_of_sign_change} & Checked & Green & Correct intermediate-value theorem argument. & None. \\
\leanname{property14_existence_under_alternative_debt_limits} & Checked & Red & The two curves are not linked to the named debt regimes. & Build both household supply curves and prove endpoint signs. \\
\leanname{assetSupply_eq_capitalDemand_at_equilibrium} & Checked & Green & Direct consequence of equilibrium clearing. & None. \\
\leanname{property15_general_equilibrium_disciplines_saving} & Checked & Green & Uses earlier equilibrium identities exactly. & None. \\
\leanname{property16_more_borrowing_reduces_capital} & Checked & Amber & A local supply comparison and excess-demand monotonicity imply the equilibrium return and capital ordering. & Derive the local supply comparison from household invariant laws. \\
\leanname{property16_more_borrowing_reduces_saving_rate} & Checked & Amber & The actual technology saving-rate function preserves the capital ordering when strictly increasing. & Derive its monotonicity from production primitives. \\
\leanname{effectiveBorrowingLimit_eq_natural_of_slack} & Checked & Green & Minimum of two ordered limits is the natural limit. & None. \\
\leanname{effectiveBorrowingLimit_admissible}, \leanname{solvedBoundedHousehold_behavior_eq_of_limit_eq}, \leanname{solvedNextExcess_eq_of_limit_eq}, \leanname{solvedResourceKernel_eq_of_limit_eq} & Checked & Green & Equal effective limits give identical solved values, policies, consumption, transitions, and kernels. & None. \\
\leanname{property17_no_effect_when_fixed_limit_slack} & Checked & Amber & Slack statutory limits now produce identical concrete household behavior, kernels, and invariant-law sets. & Build the corresponding bijection of complete equilibrium records. \\
\leanname{normalizedPath_eq} & Checked & Green & Exact growth factors cancel. & None. \\
\leanname{property18_balanced_growth_extension} & Checked & Amber & Euler wedge and growth paths are assumed. & Derive them from the growing household economy. \\
\leanname{debt_budget_translation} & Checked & Green & Exact cancellation in the debt budget. & None. \\
\leanname{debt_budget_iff_net_budget} & Checked & Green & Algebraic equivalence of gross and net budgets. & None. \\
\leanname{taxAdjusted_floor_iff_net_natural_floor} & Checked & Green & Exact translation of the borrowing floor. & None. \\
\leanname{debtAdjustedFeasible_iff_net} & Checked & Green & Pointwise budget-and-floor equivalence. & None. \\
\leanname{debtAdjustedBellmanObjective_grossAssetOfShifted}, \leanname{debtAdjustedBellmanOperator_eq_bellmanOperator}, \leanname{debtAdjustedBellmanFixedPoint_of_fixedPoint}, \leanname{debtAdjustedPolicy_optimal_of_optimal}, \leanname{debtAdjustedResourceTransition_policy} & Checked & Green & The affine coordinate change preserves the actual gross-debt objective, Bellman operator, fixed point, maximizing policy, and resource transition. & None. \\
\leanname{taxAdjustedToDebtFree_debtFreeToTaxAdjusted}, \leanname{debtFreeToTaxAdjusted_taxAdjustedToDebtFree}, \leanname{taxAdjustedAssetFloor_eq_debt_sub_natural} & Checked & Green & The complete stationary-equilibrium translations are inverse and the gross floor is exactly debt minus the natural limit. & None. \\
\leanname{fixed_floor_not_debt_neutral_witness} & Checked & Green & Explicit fixed-floor counterexample. & None. \\
\leanname{property19_debt_neutrality_depends_on_constraint} & Checked & Green & Proves a bijection of complete debt-free and tax-adjusted debt stationary equilibrium sets, plus path equivalence and the fixed-floor counterexample. & None. \\
\leanname{sharpPrice_pos} & Checked & Green & Quotient of positive terms. & None. \\
\leanname{one_div_sharpPrice_eq_naturalLimit} & Checked & Green & Exact reciprocal identity. & None. \\
\leanname{monetary_constraint_slack_above_sharp_return} & Checked & Green & Monotonic natural-limit comparison. & None. \\
\leanname{SolvedStationaryHousehold.meanAssets_eq_of_limit_eq_unique} & Checked & Green & Equal solved limits and invariant-law uniqueness identify the actual policy integrals. & None for its stated theorem. \\
\leanname{monetaryMeanAssets_eq_natural_of_slack} & Checked & Green & Slackness, solved-kernel equality, and uniqueness identify monetary and natural-limit mean assets. & None for its stated theorem. \\
\leanname{property20_monetary_return_upper_bound} & Checked & Amber & Uses concrete solved policies, invariant laws, and market-clearing zeros; only full-state uniqueness and the cross-return integral ordering remain inputs. & Derive uniqueness from Property 7 and the ordering from natural-limit asset supply. \\
\bottomrule
\end{longtable}
}

\paragraph{Interpretation.}
Green, amber, and red do not grade Lean's logical kernel: all entries are valid
theorems of their displayed hypotheses.  They grade whether those hypotheses
and conclusions constitute the foundational economic proof suggested by the
declaration's role in the Aiyagari formalization.

\end{document}
